%!TEX program = lualatex 
%!TEX encoding = UTF-8 Unicode  
%!TEX root = chapter08.tex  
\documentclass[main.tex]{subfiles} 
\begin{document} 
%----------------------------------------------------------------------------------------
%	CHAPTER 8
%----------------------------------------------------------------------------------------
\newpage 
\loadgeometry{marginlayout}  
\pagestyle{scrheadings} % Print headers again  
\KOMAoptions{
headwidth=textwithmarginpar,
headsepline=true,
headsepline= 0.5pt,
footwidth=textwithmarginpar,
}  
\onehalfspacing 
\setcounter{chapter}{8}
\chapter{Fonctions trigonométriques}  



\section{Le cercle unité}
\begin{marginfigure} \centering
\begin{tikzpicture}[scale=2.0]
\tkzSetUpPoint[size=3, fill=black]
\tkzfctset{tan style/.style={->,>=latex, color=red, line width=1.2pt}}   
\tkzSetUpLine[line width= 1.2pt, add=.5 and .5] 

\tkzInit[xmin=-1,xmax=1.0,ymin=-1,ymax=1,xstep=1, ystep=1];
%	\tkzGrid[ color=gray ](-3.0,-2.0)(6.0,5.0);   


\tkzDefPoint(0,0){O} 
%	\pgfmathsetmacro{\a}{55/360*2*pi*1}   
\tkzDefPoint(1,0){I} 
\tkzDrawCircle[line width=1.2pt, color=red](O,I)
\tkzDefPoint({cos(-1)},{sin(-1)}){T} 
\tkzText[below right](T){$x^2+y^2=1$}

%		\tkzDefPoint(-1,0){I'} 
%		\tkzDefPoint(0,1){J} 
%		\tkzDefPoint({cos(\a)},{sin(\a)}){A} 
%		\tkzDrawSegments[](O,A) 
%		\tkzLabelSegment[above left](O,A){$1$}
%		\tkzLabelPoints[below left](O){\scriptsize O}
%		\tkzLabelPoint[  right](A){\scriptsize$A(\cos(\alpha);\sin(\alpha))$}   
%		\tkzLabelAngle[pos=0.35](I,O,A){$\alpha$} 
%		\tkzMarkAngle[ arc=l,arrows=->,>=latex, size=0.2cm,mkcolor=red, line width =1.2pt](I,O,A)  
%		\tkzDefPointWith[orthogonal normed,K=1](O,A)
%		\tkzGetPoint{N} 
%\tkzDrawCircle[ line width=1.2pt](O,I) 
%		\tkzDrawArc[ line width=1.2pt](O,I)(I') 
%		\tkzPointShowCoord(A)%[ xlabel=$\cos(\alpha)$,ylabel=$\sin(\alpha)$, xstyle={below=7pt},ystyle={left=6pt}](A) 
\tkzDrawX[ noticks,font=\small, line width=2pt, label=$x$ , above =5pt ,left space=0.3,right space=0.3] 
\tkzDrawY[ noticks,font=\small, line width=2pt , label={$y$} , right =5pt, down space=0.3,up space=0.3]   
\tkzLabelX[orig=false, font=\scriptsize ]  
\tkzLabelY[orig=false,  font=\scriptsize] 
\tkzLabelPoints[below left](O)

\tkzDefPoint(0.8,0.8){A} 
\tkzLabelPoint[above right](A){$+$}
\tkzDrawArc[color=black,->,>=latex,angles,line width=2pt](O,A)(45,60)
\end{tikzpicture} 
\caption{Cercle unité}
\end{marginfigure}
\begin{definition} Le cercle trigonométrique (ou cercle unité) est le cercle de rayon 1 centré à l'origine d'équation $\mathscr{C}\colon x^2+y^2=1$.\\
Le cercle est orienté dans le sens direct (antihoraire).
\end{definition}
\begin{example} Montrer que le point $P\big(\tfrac{\sqrt{3}}{3}~;~\tfrac{\sqrt{6}}{3}\big)\in\mathscr{C}$. 
\end{example}
\begin{proof}[solution] Nous devons vérifiez que les cordonnées de $P$ vérifient l'équation du cercle unité :
\setlength{\abovedisplayskip}{3pt}
\setlength{\belowdisplayskip}{-10pt} 
$$
\left(\tfrac{\sqrt{3}}{3}\right)^2+
\left(\tfrac{\sqrt{6}}{3}\right)^2
= \tfrac{3}{9}+\tfrac{6}{9}=1
$$ 
\end{proof}



\begin{marginfigure}[1cm]\centering
\begin{tikzpicture}[scale=1.0]
\tkzSetUpPoint[size=3, fill=black]
\tkzfctset{tan style/.style={->,>=latex, color=red, line width=1.2pt}}   
\tkzSetUpLine[line width= 1.2pt, add=.5 and .5] 

\tkzInit[xmin=-2,xmax=2.0,ymin=-1.5,ymax=1.5,xstep=1, ystep=1];
%	\tkzGrid[ color=gray ](-3.0,-2.0)(6.0,5.0);   


\tkzDefPoint(0,0){O} 
%	\pgfmathsetmacro{\a}{55/360*2*pi*1}   
\tkzDefPoint(1,0){I}  
\tkzLabelPoints[below left](O)

\tkzText[above](1.25,0.75){\footnotesize Quadrant \RN{1}}
\tkzText[above](-1.25,0.75){\footnotesize Quadrant \RN{2}}
\tkzText[below](-1.25,-0.75){\footnotesize Quadrant \RN{3}}
\tkzText[below](1.25,-0.75){\footnotesize Quadrant \RN{4}}	

\tkzDrawX[noticks,above=5pt, label=$x$, font=\scriptsize, line width=1.2pt, right space=0.2pt] 
\tkzDrawY[noticks,right=5pt, label=$y$, font=\scriptsize, line width=1.2pt , up space=0.2pt,]
%	\tkzLabelX[orig=false, font=\scriptsize ]  
%	\tkzLabelY[orig=false,  font=\scriptsize] 
\end{tikzpicture} 
\caption{Les quadrants de 1 à 4}
\end{marginfigure}

\begin{definition} Dans un plan muni d'un repère orthonormé, les portions du plan délimitées par les axes du repères s'appellent \textbf{quadrants}.
\end{definition}
\begin{example} Le point $P\left(\tfrac{\sqrt{3}}{2}~;~y\right)$ appartient au cercle unité et au quadrant \RN{4}. Déterminer $y$. 
\end{example}
\begin{proof}[solution] $P\in\mathscr{C}$, $y$ vérifie $\begin{WithArrows}
\left(\tfrac{\sqrt{3}}{2}\right)^2+y^2&=1\\	
y^2& = 1-\tfrac{3}{4}=\tfrac{1}{4}\\
y& = \pm \tfrac{1}{2}
\end{WithArrows}$.

Le point $P\in$ Quadrant \RN{4}, $y$ est négatif. Donc $y=-\tfrac{1}{2}$.
\end{proof}
\begin{tBox}\textbf{L'enroulement de la droite des réels sur le cercle unité} \\
Soit $t\in\mathbb{R}$. Si $t\geqslant 0$, on place le point $P(x~;~y)$ à la distance   $t$ le long du cercle $\mathscr{C}$, en partant de $I(1~;~0)$ et en se déplaçant dans le sens \textbf{positif}. \\
Si $t<0$, on se déplace d'une distance $\abs{t}$ dans le sens \textbf{négatif}. $t$ est la mesure de l'arc orienté $\wideparen{IP}$.
\end{tBox} 
\begin{figure}[!h]
\begin{tikzpicture}[scale=2]
\tkzSetUpPoint[size=3, fill=black]%,shape=cross out
\tkzfctset{tan style/.style={->,>=latex, color=red, line width=1.2pt}}   
\tkzSetUpLine[line width= 1.2pt, add=.2 and .2]  
\tkzInit[xmin=-1.0,xmax=1.0,ymin=-1.0,ymax=1.0,xstep=1.0, ystep=1.0];  

\tkzDrawX[ font=\small, line width=2pt, label=$x$ , above =5pt ,left space=0.4,right space=0.4] 
\tkzDrawY[ font=\small, line width=2pt , label={$y$} , right =5pt, down space=0.3,up space=0.3]  
%	\tkzLabelXY[orig = false,  font=\scriptsize ] 
%	\tkzRep[color  = Maroon,xlabel=,ylabel=, font=\scriptsize, line width = 2pt,xnorm=1, ynorm=1]   
\def\a{130} 
\tkzDefPoint(0,0){O}  
\tkzDefPoint(1,0){I}  
\tkzLabelPoint[below left](O){$0$}
\tkzDrawCircle[color=blue, line width=1.2pt](O,I)
\tkzLabelPoint[below right](I){$1$}
\tkzDrawArc[color=red,->,>=latex,angles,line width=2pt](O,I)(0,\a)
\tkzDefShiftPoint[O](\a:1){P}
\tkzDrawPoints(P)
\tkzLabelPoint[above left](P){$P(x~;~y)$}	
\tkzDefShiftPoint[O]({\a/2}:1){T}	
\tkzLabelPoint[above right](T){\color{red}$t\geqslant 0$}	
%		\tkzDefShiftPoint[O](70:3){A}
%		\tkzDefShiftPoint[O](-90:3){B} 
%		\tkzDrawPoints(A,B)
%		\tkzLabelPoint[above right, font=\scriptsize](A){$A(2;y)$}
%		\tkzLabelPoint[below right, font=\scriptsize](B){$B(0;9)$}  
\end{tikzpicture} 
\begin{tikzpicture}[scale=2]
\tkzSetUpPoint[size=3, fill=black]%,shape=cross out
\tkzfctset{tan style/.style={->,>=latex, color=red, line width=1.2pt}}   
\tkzSetUpLine[line width= 1.2pt, add=.2 and .2]  
\tkzInit[xmin=-1.0,xmax=1.0,ymin=-1.0,ymax=1.0,xstep=1.0, ystep=1.0];  

\tkzDrawX[ font=\small, line width=2pt, label=$x$ , above =5pt ,left space=0.4,right space=0.4] 
\tkzDrawY[ font=\small, line width=2pt , label={$y$} , right =5pt, down space=0.3,up space=0.3]  
%	\tkzLabelXY[orig = false,  font=\scriptsize ] 
%	\tkzRep[color  = Maroon,xlabel=,ylabel=, font=\scriptsize, line width = 2pt,xnorm=1, ynorm=1]   
\def\a{-70} 
\tkzDefPoint(0,0){O}  
\tkzDefPoint(1,0){I}  
\tkzLabelPoint[below left](O){$0$}
\tkzDrawCircle[color=blue, line width=1.2pt](O,I)
\tkzLabelPoint[below right](I){$1$}
\tkzDrawArc[color=red,<-,>=latex,angles,line width=2pt](O,I)(\a,0)
\tkzDefShiftPoint[O](\a:1){P}
\tkzDrawPoints(P)
\tkzLabelPoint[below right](P){$P(x~;~y)$}	
\tkzDefShiftPoint[O]({\a/2}:1){T}	
\tkzLabelPoint[below right](T){\color{red}$t< 0$}	
%		\tkzDefShiftPoint[O](70:3){A}
%		\tkzDefShiftPoint[O](-90:3){B} 
%		\tkzDrawPoints(A,B)
%		\tkzLabelPoint[above right, font=\scriptsize](A){$A(2;y)$}
%		\tkzLabelPoint[below right, font=\scriptsize](B){$B(0;9)$}  
\end{tikzpicture} 
\caption{Point $P(x~;~y)$ du cercle unité représentant le réel $t$ dans les cas $t>0$ et $t<0$.}
\end{figure}
\noindent Le rayon du cercle unité étant $2\pi$ : 
\begin{enumerate}[label=\textbullet, leftmargin=*]
\item $0$ et $2\pi$ correspondent avec le point $P(1~;~0)$. 
\item $\tfrac{\pi}{2}=\tfrac{2\pi}{4}$ correspond avec le point $P(0~;~1)$
\end{enumerate}
\begin{figure*}[!h]
\begin{tikzpicture}[scale=1.25]
\tkzSetUpPoint[size=2, fill=black]%,shape=cross out
\tkzfctset{tan style/.style={->,>=latex, color=red, line width=1.2pt}}   
\tkzSetUpLine[line width= 1.2pt, add=.2 and .2]  
\tkzInit[xmin=-1.0,xmax=1.0,ymin=-1.0,ymax=1.0,xstep=1.0, ystep=1.0];  

\tkzDrawX[ font=\small, line width=2pt, label=$x$ , below=5pt ,left space=0.6,right space=0.6] 
\tkzDrawY[ font=\small, line width=2pt , label={$y$} , right =5pt, down space=0.5,up space=0.5]  
%	\tkzLabelXY[orig = false,  font=\scriptsize ] 
%	\tkzRep[color  = Maroon,xlabel=,ylabel=, font=\scriptsize, line width = 2pt,xnorm=1, ynorm=1]   
\def\a{90} 
\tkzDefPoint(0,0){O}  
\tkzDefPoint(1,0){I}  
\tkzLabelPoint[below left](O){$0$}
\tkzDrawCircle[color=blue, line width=1.2pt](O,I)
\tkzLabelPoint[below right](I){$1$}
\tkzDrawArc[color=red,->,>=latex,angles,line width=2pt](O,I)(0,\a)
\tkzDefShiftPoint[O](\a:1){P}
\tkzDrawPoints(P)
\tkzLabelPoint[above left](P){\footnotesize$P(0~;~1)$}	
\tkzDefShiftPoint[O]({\a/2}:1){T}	
\tkzLabelPoint[above right](T){\small\color{red}$t=\frac{\pi}{2}$}	
%		\tkzDefShiftPoint[O](70:3){A}
%		\tkzDefShiftPoint[O](-90:3){B} 
%		\tkzDrawPoints(A,B)
%		\tkzLabelPoint[above right, font=\scriptsize](A){$A(2;y)$}
%		\tkzLabelPoint[below right, font=\scriptsize](B){$B(0;9)$}  
\end{tikzpicture} 
\begin{tikzpicture}[scale=1.25]
\tkzSetUpPoint[size=2, fill=black]%,shape=cross out
\tkzfctset{tan style/.style={->,>=latex, color=red, line width=1.2pt}}   
\tkzSetUpLine[line width= 1.2pt, add=.2 and .2]  
\tkzInit[xmin=-1.0,xmax=1.0,ymin=-1.0,ymax=1.0,xstep=1.0, ystep=1.0];   
\tkzDrawX[ font=\small, line width=2pt, label=$x$ , below=5pt ,left space=0.6,right space=0.6] 
\tkzDrawY[ font=\small, line width=2pt , label={$y$} , right =5pt, down space=0.5,up space=0.5]  
%	\tkzLabelXY[orig = false,  font=\scriptsize ] 
%	\tkzRep[color  = Maroon,xlabel=,ylabel=, font=\scriptsize, line width = 2pt,xnorm=1, ynorm=1]   
\def\a{180} 
\tkzDefPoint(0,0){O}  
\tkzDefPoint(1,0){I}  
\tkzLabelPoint[below left](O){$0$}
\tkzDrawCircle[color=blue, line width=1.2pt](O,I)
\tkzLabelPoint[below right](I){$1$}
\tkzDrawArc[color=red,->,>=latex,angles,line width=2pt](O,I)(0,\a)
\tkzDefShiftPoint[O](\a:1){P}
\tkzDrawPoints(P)
\tkzLabelPoint[above left](P){\footnotesize$P(-1~;~0)$}	
\tkzDefShiftPoint[O]({60}:1){T}	
\tkzLabelPoint[above right](T){\small\color{red}$t=\pi$}	
%		\tkzDefShiftPoint[O](70:3){A}
%		\tkzDefShiftPoint[O](-90:3){B} 
%		\tkzDrawPoints(A,B)
%		\tkzLabelPoint[above right, font=\scriptsize](A){$A(2;y)$}
%		\tkzLabelPoint[below right, font=\scriptsize](B){$B(0;9)$}  
\end{tikzpicture}  
\begin{tikzpicture}[scale=1.25]
\tkzSetUpPoint[size=2, fill=black]%,shape=cross out
\tkzfctset{tan style/.style={->,>=latex, color=red, line width=1.2pt}}   
\tkzSetUpLine[line width= 1.2pt, add=.2 and .2]  
\tkzInit[xmin=-1.0,xmax=1.0,ymin=-1.0,ymax=1.0,xstep=1.0, ystep=1.0];   
\tkzDrawX[ font=\small, line width=2pt, label=$x$ , below=5pt ,left space=0.6,right space=0.6] 
\tkzDrawY[ font=\small, line width=2pt , label={$y$} , right =5pt, down space=0.5,up space=0.5]  
%	\tkzLabelXY[orig = false,  font=\scriptsize ] 
%	\tkzRep[color  = Maroon,xlabel=,ylabel=, font=\scriptsize, line width = 2pt,xnorm=1, ynorm=1]   
\def\a{270} 
\tkzDefPoint(0,0){O}  
\tkzDefPoint(1,0){I}  
\tkzLabelPoint[below left](O){$0$}
\tkzDrawCircle[color=blue, line width=1.2pt](O,I)
\tkzLabelPoint[below right](I){$1$}
\tkzDrawArc[color=red,->,>=latex,angles,line width=2pt](O,I)(0,\a)
\tkzDefShiftPoint[O](\a:1){P}
\tkzDrawPoints(P)
\tkzLabelPoint[below left](P){\footnotesize$P(0~;~-1)$}	
\tkzDefShiftPoint[O]({60}:1){T}	
\tkzLabelPoint[above right](T){\small\color{red}$t=\frac{3\pi}{2}$}	
%		\tkzDefShiftPoint[O](70:3){A}
%		\tkzDefShiftPoint[O](-90:3){B} 
%		\tkzDrawPoints(A,B)
%		\tkzLabelPoint[above right, font=\scriptsize](A){$A(2;y)$}
%		\tkzLabelPoint[below right, font=\scriptsize](B){$B(0;9)$}  
\end{tikzpicture}  
\begin{tikzpicture}[scale=1.25]
\tkzSetUpPoint[size=2, fill=black]%,shape=cross out
\tkzfctset{tan style/.style={->,>=latex, color=red, line width=1.2pt}}   
\tkzSetUpLine[line width= 1.2pt, add=.2 and .2]  
\tkzInit[xmin=-1.0,xmax=1.0,ymin=-1.0,ymax=1.0,xstep=1.0, ystep=1.0];   
\tkzDrawX[ font=\small, line width=2pt, label=$x$ , below=5pt ,left space=0.6,right space=0.6] 
\tkzDrawY[ font=\small, line width=2pt , label={$y$} , right =5pt, down space=0.5,up space=0.5]  
%	\tkzLabelXY[orig = false,  font=\scriptsize ] 
%	\tkzRep[color  = Maroon,xlabel=,ylabel=, font=\scriptsize, line width = 2pt,xnorm=1, ynorm=1]   
\def\a{360} 
\tkzDefPoint(0,0){O}  
\tkzDefPoint(1,0){I}  
\tkzLabelPoint[below left](O){$0$}
\tkzDrawCircle[color=blue, line width=1.2pt](O,I)
\tkzLabelPoint[below right](I){$1$}
\tkzDrawArc[color=red,->,>=latex,angles,line width=2pt](O,I)(0,\a)
\tkzDefShiftPoint[O](\a:1){P}
\tkzDrawPoints(P)
\tkzLabelPoint[above right](P){\footnotesize$P(1~;~0)$}	
\tkzDefShiftPoint[O]({60}:1){T}	
\tkzLabelPoint[above right](T){\small\color{red}$t=2\pi$}	
%		\tkzDefShiftPoint[O](70:3){A}
%		\tkzDefShiftPoint[O](-90:3){B} 
%		\tkzDrawPoints(A,B)
%		\tkzLabelPoint[above right, font=\scriptsize](A){$A(2;y)$}
%		\tkzLabelPoint[below right, font=\scriptsize](B){$B(0;9)$}  
\end{tikzpicture}   
\caption{Points associés aux réels $t=\tfrac{\pi}{2} =\tfrac{2\pi}{4}$ quart de tour,~$\pi =\tfrac{2\pi}{2}$ demi tour,~$\tfrac{3\pi}{2}$ (3 quarts de tour) et $2\pi$ (tour complet).}
\end{figure*}  
\begin{figure*}[!h]
\begin{tikzpicture}[scale=1.25]
\tkzSetUpPoint[size=2, fill=black]%,shape=cross out
\tkzfctset{tan style/.style={->,>=latex, color=red, line width=1.2pt}}   
\tkzSetUpLine[line width= 1.2pt, add=.2 and .2]  
\tkzInit[xmin=-1.35,xmax=1.35,ymin=-1.25,ymax=1.25,xstep=1.0, ystep=1.0];  

\tkzDrawX[ noticks, font=\small, line width=2pt, label=$x$ , below=5pt ,left space=0.5,right space=0.5] 
\tkzDrawY[ noticks,  font=\small, line width=2pt , label={$y$} , right =5pt, down space=0.4,up space=0.4]  
%	\tkzLabelXY[orig = false,  font=\scriptsize ] 
%	\tkzRep[color  = Maroon,xlabel=,ylabel=, font=\scriptsize, line width = 2pt,xnorm=1, ynorm=1]   
\def\a{540}  
\tkzDefPoint(0,0){O}  
\tkzDefPoint(1,0){I}  
\tkzLabelPoint[below left](O){$0$}
\tkzDrawCircle[color=blue, line width=1.2pt](O,I)
\tkzLabelPoint[below right](I){$1$}
\tkzFctPar[smooth,samples=300,domain=0:9.425,->,>=latex,line width=0.8pt,color=red]{(1+3*t/94.25)*cos(t)}{(1+3*t/94.25)*sin(t)}
%		\tkzDrawArc[color=red,->,>=latex,angles,line width=2pt](O,I)(0,\a) 
\tkzDefShiftPoint[O](\a:1){P}
\tkzDrawPoints(P)
\tkzLabelPoint[above left](P){\footnotesize$P(-1~;~0)$}	
\tkzDefShiftPoint[O]({60}:1){T}	
\tkzLabelPoint[above right](T){\small\color{red}$t=3\pi$}	
%		\tkzDefShiftPoint[O](70:3){A}
%		\tkzDefShiftPoint[O](-90:3){B} 
%		\tkzDrawPoints(A,B)
%		\tkzLabelPoint[above right, font=\scriptsize](A){$A(2;y)$}
%		\tkzLabelPoint[below right, font=\scriptsize](B){$B(0;9)$}  
\end{tikzpicture}  
\begin{tikzpicture}[scale=1.25]
\tkzSetUpPoint[size=2, fill=black]%,shape=cross out
\tkzfctset{tan style/.style={->,>=latex, color=red, line width=1.2pt}}   
\tkzSetUpLine[line width= 1.2pt, add=.2 and .2]  
\tkzInit[xmin=-1.250,xmax=1.25,ymin=-1.25,ymax=1.25,xstep=1.0, ystep=1.0];  
\tkzDrawX[ noticks, font=\small, line width=2pt, label=$x$ , below=5pt ,left space=0.6,right space=0.6] 
\tkzDrawY[ noticks,  font=\small, line width=2pt , label={$y$} , right =5pt, down space=0.4,up space=0.4]  
%	\tkzLabelXY[orig = false,  font=\scriptsize ] 
%	\tkzRep[color  = Maroon,xlabel=,ylabel=, font=\scriptsize, line width = 2pt,xnorm=1, ynorm=1]   
\def\a{180} 
\tkzDefPoint(0,0){O}  
\tkzDefPoint(1,0){I}  
\tkzLabelPoint[below left](O){$0$}
\tkzDrawCircle[color=blue, line width=1.2pt](O,I)
\tkzLabelPoint[below right](I){$1$}
\tkzDrawArc[color=red,<-,>=latex,angles,line width=2pt](O,I)(\a,0)
\tkzDefShiftPoint[O](\a:1){P}
\tkzDrawPoints(P)
\tkzLabelPoint[above left](P){\footnotesize$P(-1~;~0)$}	
\tkzDefShiftPoint[O]({-60}:1){T}	
\tkzLabelPoint[below right](T){\small\color{red}$t=-\pi$}	
%		\tkzDefShiftPoint[O](70:3){A}
%		\tkzDefShiftPoint[O](-90:3){B} 
%		\tkzDrawPoints(A,B)
%		\tkzLabelPoint[above right, font=\scriptsize](A){$A(2;y)$}
%		\tkzLabelPoint[below right, font=\scriptsize](B){$B(0;9)$}  
\end{tikzpicture}  
\begin{tikzpicture}[scale=1.25]
\tkzSetUpPoint[size=2, fill=black]%,shape=cross out
\tkzfctset{tan style/.style={->,>=latex, color=red, line width=1.2pt}}   
\tkzSetUpLine[line width= 1.2pt, add=.2 and .2]  
\tkzInit[xmin=-1.250,xmax=1.25,ymin=-1.25,ymax=1.25,xstep=1.0, ystep=1.0];  
\tkzDrawX[ noticks, font=\small, line width=2pt, label=$x$ , below=5pt ,left space=0.6,right space=0.6] 
\tkzDrawY[ noticks,  font=\small, line width=2pt , label={$y$} , right =5pt, down space=0.4,up space=0.4]  
%	\tkzLabelXY[orig = false,  font=\scriptsize ] 
%	\tkzRep[color  = Maroon,xlabel=,ylabel=, font=\scriptsize, line width = 2pt,xnorm=1, ynorm=1]   
\def\a{-90} 
\tkzDefPoint(0,0){O}  
\tkzDefPoint(1,0){I}  
\tkzLabelPoint[below left](O){$0$}
\tkzDrawCircle[color=blue, line width=1.2pt](O,I)
\tkzLabelPoint[below right](I){$1$}
\tkzDrawArc[color=red,<-,>=latex,angles,line width=2pt](O,I)(\a,0)
\tkzDefShiftPoint[O](\a:1){P}
\tkzDrawPoints(P)
\tkzLabelPoint[below left](P){\footnotesize$P(0~;~-1)$}	
\tkzDefShiftPoint[O]({-60}:1){T}	
\tkzLabelPoint[below  right](T){\small\color{red}$t=-\frac{\pi}{2}$}	
%		\tkzDefShiftPoint[O](70:3){A}
%		\tkzDefShiftPoint[O](-90:3){B} 
%		\tkzDrawPoints(A,B)
%		\tkzLabelPoint[above right, font=\scriptsize](A){$A(2;y)$}
%		\tkzLabelPoint[below right, font=\scriptsize](B){$B(0;9)$}  
\end{tikzpicture}   
\caption{Les points correspondant aux réels $t=3\pi$, ~$-\pi$ et $-\tfrac{\pi}{2}$.}
\end{figure*} 
\begin{example} Donner les coordonnées du point associé à $t=\tfrac{\pi}{4}$.
\begin{marginfigure} 
\begin{tikzpicture}[scale=1.25]
\tkzSetUpPoint[size=2, fill=black]%,shape=cross out
\tkzfctset{tan style/.style={->,>=latex, color=red, line width=1.2pt}}   
\tkzSetUpLine[line width= 1.2pt, add=.2 and .2]  
\tkzInit[xmin=-1.250,xmax=1.25,ymin=-1.25,ymax=1.25,xstep=1.0, ystep=1.0];  
\tkzDrawX[ noticks, font=\small, line width=2pt, label=$x$ , below=5pt ,left space=0.6,right space=0.6] 
\tkzDrawY[ noticks,  font=\small, line width=2pt , label={$y$} , right =5pt, down space=0.2,up space=0.3]  
%	\tkzLabelXY[orig = false,  font=\scriptsize ] 
%	\tkzRep[color  = Maroon,xlabel=,ylabel=, font=\scriptsize, line width = 2pt,xnorm=1, ynorm=1]   
\def\a{45} 
\tkzDefPoint(0,0){O}  
\tkzDefPoint(1,0){I}  
\tkzLabelPoint[below left](O){$0$}
\tkzDrawCircle[color=blue, line width=1.2pt](O,I)
\tkzLabelPoint[below right](I){$1$}
\tkzDrawArc[color=red,->,>=latex,angles,line width=2pt](O,I)(0,\a)
\tkzDefShiftPoint[O](\a:1){P}
\tkzDrawPoints(P)
\tkzFct[samples=2,line width=1.2pt,]{\x}
%	\tkzLabelPoint[below left](P){\footnotesize$P(\tfrac{\sqrt{2}}{2}~;~\tfrac{\sqrt{2}}{2})$}	
\tkzDefShiftPoint[O]({\a/2}:1){T}	
\tkzLabelPoint[  right](T){\small\color{red}$t=\frac{\pi}{4}$}	
%		\tkzDefShiftPoint[O](70:3){A}
%		\tkzDefShiftPoint[O](-90:3){B} 
%		\tkzDrawPoints(A,B)
%		\tkzLabelPoint[above right, font=\scriptsize](A){$A(2;y)$}
%		\tkzLabelPoint[below right, font=\scriptsize](B){$B(0;9)$}  
\end{tikzpicture}   
\caption{$t=\frac{\pi}{4}$}
\end{marginfigure}  
\end{example}
\begin{proof}[solution] $P(x~;~y)$ est sur la droite d'équation $y=x$.\\
$x$ vérifie $\begin{WithArrows}
x^2+x^2&=1\\	
2x^2& = 1\\
x& = \pm \tfrac{\sqrt{2}}{2}
\end{WithArrows}$  \\
$P$ est dans le quadrant \RN{1}, $x$ et $y$ sont positifs et $x=y=\tfrac{\sqrt{2}}{2}$
\end{proof}
\begin{table*}[!h] 
\renewcommand{\arraystretch}{1.8}   \begin{tabularx}{1\linewidth}{|>{\columncolor{blue!20}}c|Y|Y|Y|Y|Y|} \hline
\textbf{t}& $0$ & $\tfrac{\pi}{6}$ &$\tfrac{\pi}{4}$&$\tfrac{\pi}{3}$&$\tfrac{\pi}{2}$ \\ \hline 
point $P$ & 	$\left(1~;~0\right)$
& 	$\left(\frac{\sqrt{3}}{2}~;~\frac{1}{2}\right)$
& 	$\left(\frac{\sqrt{2}}{2}~;~\frac{\sqrt{2}}{2}\right)$
	& 	$\left(\frac{1}{2}~;~\frac{\sqrt{3}}{2}\right)$
& 	$\left(0~;~1\right)$ \\ \hline
\end{tabularx}
\caption{Valeurs à retenir \emoji{heart}}
\end{table*}
%----------------------------------------------------------------------------------------
%	 
%----------------------------------------------------------------------------------------
\newpage 
\loadgeometry{widelayout}  
\pagestyle{scrheadings} % Print headers again  
\KOMAoptions{
headwidth=textwithmarginpar,
headsepline=true,
headsepline= 0.5pt,
footwidth=textwithmarginpar,
}  
\onehalfspacing 
\subsection{Exercices cercle trigonométrique}

\begin{exercise}[concepts] Complétez
\begin{spacing}{2}
\begin{enumerate}[label={},leftmargin=*]
\item Le cercle trigonométrique est un cercle de centre \dotfill et de rayon \dotfill
\item Dans un repère orthonormé, l'équation du cercle trigonométrique est \dotfill 
\item Les points $A(1~;~\ldots\ldots)$, $B(-1~;~\ldots\ldots)$,$C(\ldots\ldots~;~1)$ et $D(\ldots\ldots~;~-1)$ sont sur $\mathscr{C}$. 
\item Les points correspondants aux réels $\tfrac{\pi}{2}$; $\pi$; $-\tfrac{\pi}{2}$ et $2\pi$ ont respectivement pour coordonnées \dotfill, \\
\null \dotfill,\dotfill et \dotfill.
\end{enumerate}
\end{spacing} \vspace{-1em}
\end{exercise}  
\begin{exercise} Montrer que le point donné est sur le cercle unité.
\vspace{-0.5em}
\begin{spacing}{1.8}
\begin{multicols}{4} 
\begin{enumerate}[label={\arabic*)},leftmargin=*]
\item $A\left(\tfrac{3}{5}~;~-\tfrac{4}{5}\right)$
\item $B\left(\tfrac{-24}{25}~;~-\tfrac{7}{25}\right)$
\item $C\left(\tfrac{3}{4}~;~-\tfrac{\sqrt{7}}{4}\right)$
\item $D\left(-\tfrac{5}{7}~;~-\tfrac{2\sqrt{6}}{7}\right)$
\end{enumerate}
\end{multicols}\vspace{-2em}
\end{spacing}
\end{exercise}
\begin{exercise} Déterminer la coordonnée manquante sachant que $P\in\mathscr{C}$ et le quadrant indiqué.
\vspace{-0.5em}
\begin{spacing}{1.8}
\begin{multicols}{3} 
\begin{enumerate}[label={\arabic*)},leftmargin=*]
\item $P(-\tfrac{3}{5}~;~\ldots) \in$ Quadrant \RN{3}
\item  $P(\ldots~;~-\tfrac{7}{25}) \in$ Quadrant \RN{4}
\item  $P(\ldots~;~\tfrac{1}{3})  \in$ Quadrant \RN{2}
%		\item  $P(\tfrac{2}{5}~;~\ldots)$ et $P\in$ Quadrant \RN{1}
%		\item  $P(\ldots~;~-\tfrac{2}{7})$ et $P\in$ Quadrant \RN{4}
%		\item  $P(-\tfrac{2}{3}~;~\ldots)$ et $P\in$ Quadrant \RN{2}
\end{enumerate}
\end{multicols}\vspace{-2em}
\end{spacing}	
\end{exercise}
\begin{exercise} Les cercles trigonométriques sont marqués avec $t$ augmentant par  incréments de $\tfrac{\pi}{4}$ et $\tfrac{\pi}{6}$ respectivement. Complétez les coordonnées des points indqués.\\
\noindent\parbox{0.48\linewidth}{\centering
\begin{tikzpicture}[scale=2.8]
\tkzSetUpPoint[size=4, fill=black]%,shape=cross out
\tkzfctset{tan style/.style={->,>=latex, color=red, line width=1.2pt}}   
\tkzSetUpLine[line width= 1.2pt, add=.2 and .2]  
\tkzInit[xmin=-1.0,xmax=1.0,ymin=-1.0,ymax=1.0,xstep=1.0, ystep=1.0];   
\tkzDefPoint(0,0){O}  
\tkzDefPoint(1,0){I}   
\tkzDefPoint(0,1){J}  
\tkzDrawX[ font=\small, line width=1.2pt, label={}  ,below left, left space=0.1,right space=0.1] 
\tkzDrawY[ font=\small, line width=1.2pt , label={} ,below left, down space=0.1,up space=0.1]  
\tkzDrawCircle[line width=2pt,color=black](O,I)
\foreach \x/\t/\xtext/\ytext in {
0/0/1/0,
%		30/\tfrac{\pi}{6}/\tfrac{\sqrt{3}}{2}/\tfrac{1}{2},
%		60/\tfrac{\pi}{3}/\tfrac{1}{2}/\tfrac{\sqrt{3}}{2},
45/\tfrac{\pi}{4}/\tfrac{\sqrt{2}}{2}/\tfrac{\sqrt{2}}{2}%,
%		90/\tfrac{\pi}{2}/0/1, 
%		135/\tfrac{3\pi}{4}/-\tfrac{\sqrt{2}}{2}/\tfrac{\sqrt{2}}{2}, 
%		180/\pi/-1/0,
%		-45/-\tfrac{\pi}{4}/\tfrac{\sqrt{2}}{2}/-\tfrac{\sqrt{2}}{2},
%		-90/\tfrac{\pi}{2}/0/1
%		-135/-\tfrac{3\pi}{4}/-\tfrac{\sqrt{2}}{2}/-\tfrac{\sqrt{2}}{2}
} 
{
\tkzDefShiftPoint[O](\x:1){P};
\tkzDrawPoints(P);
\tkzDefShiftPoint[O](\x:0.95){S};
\tkzDefShiftPoint[O](\x:1.05){T};
\tkzDrawSegment[line width=1.2pt](S,T);
\tkzDrawSegment[line width=1pt](O,P);
%		\tkzText[fill=white](\x:0.60){\scriptsize\ang{\x}};
%		\tkzText[fill=white](\x:1.2){\scriptsize$t=\t$};
\tkzText[text width=1.5cm](\x:1.45){\scriptsize$P\left(\xtext,\ytext\right)$\newline \scriptsize$t=\t$};
}
\foreach \x/\t/\xtext/\ytext in { 
	90/\tfrac{\pi}{2}/0/1, 
	135/\tfrac{3\pi}{4}/-\tfrac{\sqrt{2}}{2}/\tfrac{\sqrt{2}}{2}, 
	180/\pi/-1/0,
	-45/-\tfrac{\pi}{4}/\tfrac{\sqrt{2}}{2}/-\tfrac{\sqrt{2}}{2},
	-90/\tfrac{\pi}{2}/0/1,
	-135/-\tfrac{3\pi}{4}/-\tfrac{\sqrt{2}}{2}/-\tfrac{\sqrt{2}}{2}
} 
{
\tkzDefShiftPoint[O](\x:1){P};
\tkzDrawPoints(P);
\tkzDefShiftPoint[O](\x:0.95){S};
\tkzDefShiftPoint[O](\x:1.05){T};
\tkzDrawSegment[line width=1.2pt](S,T);
\tkzDrawSegment[line width=1pt](O,P);
%		\tkzText[fill=white](\x:0.60){\scriptsize\ang{\x}};
%	\tkzText[fill=white](\x:1.25){$t=\qquad$};
\tkzText[text width=1.5cm]((\x:1.45){$P\left(\quad~;~\quad\right)$\newline $t=\qquad$};
} 
\end{tikzpicture} 
}
%
\hfill
%\noindent
\parbox{0.5\linewidth}{\centering
\begin{tikzpicture}[scale=3.0]
\tkzSetUpPoint[size=4, fill=black]%,shape=cross out
\tkzfctset{tan style/.style={->,>=latex, color=red, line width=1.2pt}}   
\tkzSetUpLine[line width= 1.2pt, add=.2 and .2]  
\tkzInit[xmin=-1.0,xmax=1.0,ymin=-1.0,ymax=1.0,xstep=1.0, ystep=1.0];   
\tkzDefPoint(0,0){O}  
\tkzDefPoint(1,0){I}   
\tkzDefPoint(0,1){J}  
\tkzDrawX[ font=\small, line width=1.2pt, label={}  ,below left, left space=0.1,right space=0.1] 
\tkzDrawY[ font=\small, line width=1.2pt , label={} ,below left, down space=0.1,up space=0.1]  
\tkzDrawCircle[line width=2pt,color=black](O,I)
\foreach \x/\t/\xtext/\ytext in {
0/0/1/0,
30/\tfrac{\pi}{6}/\tfrac{\sqrt{3}}{2}/\tfrac{1}{2}
%		60/\tfrac{\pi}{3}/\tfrac{1}{2}/\tfrac{\sqrt{3}}{2},
%		45/\tfrac{\pi}{4}/\tfrac{\sqrt{2}}{2}/\tfrac{\sqrt{2}}{2}%,
%		90/\tfrac{\pi}{2}/0/1, 
%		135/\tfrac{3\pi}{4}/-\tfrac{\sqrt{2}}{2}/\tfrac{\sqrt{2}}{2}, 
%		180/\pi/-1/0,
%		-45/-\tfrac{\pi}{4}/\tfrac{\sqrt{2}}{2}/-\tfrac{\sqrt{2}}{2},
%		-90/\tfrac{\pi}{2}/0/1
%		-135/-\tfrac{3\pi}{4}/-\tfrac{\sqrt{2}}{2}/-\tfrac{\sqrt{2}}{2}
} 
{
\tkzDefShiftPoint[O](\x:1){P};
\tkzDrawPoints(P);
\tkzDefShiftPoint[O](\x:0.95){S};
\tkzDefShiftPoint[O](\x:1.05){T};
\tkzDrawSegment[line width=1.2pt](S,T);
\tkzDrawSegment[line width=1pt](O,P);
%		\tkzText[fill=white](\x:0.60){\scriptsize\ang{\x}};
%		\tkzText[fill=white](\x:1.2){\scriptsize$t=\t$};
\tkzText[text width=1.5cm](\x:1.45){\scriptsize$P\left(\xtext,\ytext\right)$\newline \scriptsize$t=\t$};
}
\foreach \x/\t/\xtext/\ytext in { 
	60/\tfrac{\pi}{3}/\tfrac{1}{2}/\tfrac{\sqrt{3}}{2}, 
%		90/\tfrac{\pi}{2}/0/1,  
%		180/\pi/-1/0, 
%		-90/\tfrac{\pi}{2}/0/1,
120/\tfrac{2\pi}{3}/-\tfrac{1}{2}/\tfrac{\sqrt{3}}{2},
150/\tfrac{5\pi}{6}/-\tfrac{\sqrt{3}}{2}/\tfrac{1}{2},
-30/-\tfrac{\pi}{6}/\tfrac{\sqrt{3}}{2}/-\tfrac{1}{2},
-60/-\tfrac{\pi}{3}/\tfrac{1}{2}/-\tfrac{\sqrt{3}}{2},
-120/-\tfrac{2\pi}{3}/-\tfrac{1}{2}/-\tfrac{\sqrt{3}}{2},
-150/-\tfrac{5\pi}{6}/-\tfrac{\sqrt{3}}{2}/-\tfrac{1}{2} 
} 
{
\tkzDefShiftPoint[O](\x:1){P};
\tkzDrawPoints(P);
\tkzDefShiftPoint[O](\x:0.95){S};
\tkzDefShiftPoint[O](\x:1.05){T};
\tkzDrawSegment[line width=1.2pt](S,T);
\tkzDrawSegment[line width=1pt](O,P);
%		\tkzText[fill=white](\x:0.60){\scriptsize\ang{\x}};
%	\tkzText[fill=white](\x:1.25){$t=\qquad$};
\tkzText[text width=1.5cm]((\x:1.5){$P\left(\quad~;~\quad\right)$\newline $t=\qquad$};
}
%	\foreach \x/\t/\xtext/\ytext in { 
%		120/\tfrac{2\pi}{3}/-\tfrac{1}{2}/\tfrac{\sqrt{3}}{2},
%		150/\tfrac{5\pi}{6}/-\tfrac{\sqrt{3}}{2}/\tfrac{1}{2},
%		-30/-\tfrac{\pi}{6}/\tfrac{\sqrt{3}}{2}/-\tfrac{1}{2},
%		-60/-\tfrac{\pi}{3}/\tfrac{1}{2}/-\tfrac{\sqrt{3}}{2},
%		-120/-\tfrac{2\pi}{3}/-\tfrac{1}{2}/-\tfrac{\sqrt{3}}{2},
%		-150/-\tfrac{5\pi}{6}/-\tfrac{\sqrt{3}}{2}/-\tfrac{1}{2},
%		-90/\tfrac{\pi}{2}/0/1
%	} 
%	{
%		\tkzDefShiftPoint[O](\x:1){P};
%		\tkzDefShiftPoint[O](\x:0.95){S};
%		\tkzDefShiftPoint[O](\x:1.05){T};
%		\tkzDrawSegment[line width=1.2pt](S,T);
%		\tkzDefShiftPoint[O](\x:1){T};
%		\tkzDrawPoints(P);
%		%		\tkzDrawSegment[line width=1pt](O,P);
%		%		\tkzText[fill=white](\x:0.60){\scriptsize\ang{\x}};
%		\tkzText[fill=white](\x:1.15){\scriptsize$\t$};
%		\tkzText(\x:1.40){\scriptsize$\left(\xtext,\ytext\right)$};
%	}
\end{tikzpicture} 
} 
\end{exercise}
\begin{exercise}[modèle] Déterminer les coordonnées des points correspondants aux réels.
\vspace{-0.75em}
\begin{spacing}{1.8}
\begin{multicols}{3} 
	\begin{enumerate}[label={\alph*)},leftmargin=*]
		\item $t=-\tfrac{\pi}{4}$
		\item $t=\tfrac{3\pi}{4}$
		\item $t=-\tfrac{5\pi}{6}$
	\end{enumerate}
\end{multicols}\vspace{-2.75em}
\end{spacing}	
\end{exercise}
\begin{proof}[solution] Complétez :\vspace{-1em}\\ 
\noindent
\parbox{0.72\linewidth}{ 
%	\vspace{-0.5em}
\begin{spacing}{1.8}
\begin{enumerate}[label=\alph*),leftmargin=*] 
\item $P$ le point correspondant à $-\tfrac{\pi}{4}$, et $Q$ le point correspondant à $\tfrac{\pi}{4}$.\\
( placer $P$ et $Q$ sur le cercle trigonométrique )\\
Les points sont symétriques par rapport \dotfill  \\
Comme $Q\left(\tfrac{\sqrt{2}}{2}~;~\tfrac{\sqrt{2}}{2}\right)$, on a donc $P$\dotfill 
\item $P$ le point correspondant à $\tfrac{3\pi}{4}$, et $Q$ le point correspondant à $\tfrac{\pi}{4}$.\\
	( placer $P$ et $Q$ sur le cercle trigonométrique )\\
	Les points sont symétriques par rapport \dotfill  \\
	Comme $Q\left(\tfrac{\sqrt{2}}{2}~;~\tfrac{\sqrt{2}}{2}\right)$, on a donc $P$\dotfill 
\item $P$ le point correspondant à $-\tfrac{5\pi}{6}$, et $Q$ le point correspondant à $\tfrac{\pi}{6}$.\\
	( placer $P$ et $Q$ sur le cercle trigonométrique )\\
	Les points sont symétriques par rapport \dotfill  \\
	Comme $Q\left(\qquad ~;~\qquad \right)$, on a donc $P$\dotfill 
\end{enumerate}\vspace{-2.5em}
\end{spacing}
}
\hfill
\parbox{0.22\linewidth}{
	\begin{tikzpicture}[scale=1.5]
	\tkzSetUpPoint[size=3, fill=black]%,shape=cross out
	\tkzfctset{tan style/.style={->,>=latex, color=red, line width=1.2pt}}   
	\tkzSetUpLine[line width= 1.2pt, add=.2 and .2]  
	\tkzInit[xmin=-1.0,xmax=1.0,ymin=-1.0,ymax=1.0,xstep=1.0, ystep=1.0];  
	
	%	\tkzLabelXY[orig = false,  font=\scriptsize ] 
	%	\tkzRep[color  = Maroon,xlabel=,ylabel=, font=\scriptsize, line width = 2pt,xnorm=1, ynorm=1]   
	\def\q{45} 
	\def\p{-45} 
	\tkzDefPoint(0,0){O}  
	\tkzDefPoint(1,0){I}  
	\tkzDefPoint(0,1){J}
	\tkzLabelPoint[below left](O){$0$}
	\tkzDrawCircle[color=blue, line width=1.2pt](O,I)
	\tkzLabelPoint[below right](I){$1$}
	\tkzLabelPoint[above right](J){$1$}
	\tkzDrawArc[color=red,->,>=latex,angles,line width=2pt](O,I)(0,\q)
	\tkzDrawArc[color=red,<-,>=latex,angles,line width=2pt](O,I)(\p,0)
	\tkzDefShiftPoint[O](\q:1){Q}
	\tkzDrawPoints(Q)
	\tkzLabelPoint[above right](Q){$Q$}	
	\tkzDefShiftPoint[O](\p:1){P}
	\tkzDrawPoints(P)
	\tkzLabelPoint[below right](P){$P$}	  
	 
	\tkzDrawX[ font=\small, line width=2pt, label=$x$  ,left space=0.1,right space=0.4] 
	\tkzDrawY[ font=\small, line width=2pt , label={$y$} , down space=0.1,up space=0.3]  
\end{tikzpicture} \\
	\begin{tikzpicture}[scale=1.5]
	\tkzSetUpPoint[size=3, fill=black]%,shape=cross out
	\tkzfctset{tan style/.style={->,>=latex, color=red, line width=1.2pt}}   
	\tkzSetUpLine[line width= 1.2pt, add=.2 and .2]  
	\tkzInit[xmin=-1.0,xmax=1.0,ymin=-1.0,ymax=1.0,xstep=1.0, ystep=1.0];  
	
	%	\tkzLabelXY[orig = false,  font=\scriptsize ] 
	%	\tkzRep[color  = Maroon,xlabel=,ylabel=, font=\scriptsize, line width = 2pt,xnorm=1, ynorm=1]   
	\def\q{45} 
	\def\p{3*45} 
	\tkzDefPoint(0,0){O}  
	\tkzDefPoint(1,0){I}  
	\tkzDefPoint(0,1){J}
	\tkzLabelPoint[below left](O){$0$}
	\tkzDrawCircle[color=blue, line width=1.2pt](O,I)
	\tkzLabelPoint[below right](I){$1$}
	\tkzLabelPoint[above right](J){$1$}
%	\tkzDrawArc[color=red,->,>=latex,angles,line width=2pt](O,I)(0,\q)
	\tkzDrawArc[color=red,->,>=latex,angles,line width=2pt](O,I)(0,\p)
%	\tkzDefShiftPoint[O](\q:1){Q}
%	\tkzDrawPoints(Q)
%	\tkzLabelPoint[above right](Q){$Q$}	
	\tkzDefShiftPoint[O](\p:1){P}
	\tkzDrawPoints(P)
	\tkzLabelPoint[above left](P){$P$}	  
	
	\tkzDrawX[ font=\small, line width=2pt, label=$x$  ,left space=0.1,right space=0.4] 
	\tkzDrawY[ font=\small, line width=2pt , label={$y$} , down space=0.1,up space=0.3]  
\end{tikzpicture} \\
\begin{tikzpicture}[scale=1.5]
\tkzSetUpPoint[size=3, fill=black]%,shape=cross out
\tkzfctset{tan style/.style={->,>=latex, color=red, line width=1.2pt}}   
\tkzSetUpLine[line width= 1.2pt, add=.2 and .2]  
\tkzInit[xmin=-1.0,xmax=1.0,ymin=-1.0,ymax=1.0,xstep=1.0, ystep=1.0];  

%	\tkzLabelXY[orig = false,  font=\scriptsize ] 
%	\tkzRep[color  = Maroon,xlabel=,ylabel=, font=\scriptsize, line width = 2pt,xnorm=1, ynorm=1]   
\def\q{30} 
\def\p{-5*30} 
\tkzDefPoint(0,0){O}  
\tkzDefPoint(1,0){I}  
\tkzDefPoint(0,1){J}
\tkzLabelPoint[below left](O){$0$}
\tkzDrawCircle[color=blue, line width=1.2pt](O,I)
\tkzLabelPoint[below right](I){$1$}
\tkzLabelPoint[above right](J){$1$}
%\tkzDrawArc[color=red,->,>=latex,angles,line width=2pt](O,I)(0,\q)
\tkzDrawArc[color=red,<-,>=latex,angles,line width=2pt](O,I)(\p,0)
%\tkzDefShiftPoint[O](\q:1){Q}
%\tkzDrawPoints(Q)
%\tkzLabelPoint[above right](Q){$Q$}	
\tkzDefShiftPoint[O](\p:1){P}
\tkzDrawPoints(P)
\tkzLabelPoint[below left](P){$P$}	  

\tkzDrawX[ font=\small, line width=2pt, label=$x$  ,left space=0.1,right space=0.4] 
\tkzDrawY[ font=\small, line width=2pt , label={$y$} , down space=0.1,up space=0.3]  
\end{tikzpicture} 
}
\end{proof}
\begin{exercise}Pour chaque $t\in\mathbb{R}$, tracer un cercle trigonométrique à main levée et retrouver les coordonnées du point associé  :
	\vspace{-0.75em}
	\begin{spacing}{1.8}
		\begin{multicols}{5} 
			\begin{enumerate}[label={\arabic*)},leftmargin=*]
				\item $t=4\pi$
				\item $t=-\tfrac{\pi}{6}$ 
				\item $t=\tfrac{3\pi}{2}$ 
				\item $t=-3\pi$  
				\item $t=\tfrac{4\pi}{3}$ 
%				\item $t=\tfrac{5\pi}{2}$ 
%				\item $t=\tfrac{5\pi}{4}$ 
%				\item $t=-\tfrac{7\pi}{4}$ 
%				\item $t=\tfrac{11\pi}{6}$ 
%				\item $t=\tfrac{5\pi}{3}$ 
			\end{enumerate}
		\end{multicols}\vspace{-2.75em}
	\end{spacing}	
\end{exercise}
\begin{exercise}[entrainement] Mêmes consignes
	\vspace{-0.75em}
	\begin{spacing}{1.8}
		\begin{multicols}{5} 
			\begin{enumerate}[label={\arabic*)},leftmargin=*]
				\item $t=\tfrac{5\pi}{2}$ 
				\item $t=\tfrac{5\pi}{4}$ 
				\item $t=-\tfrac{7\pi}{4}$ 
				\item $t=\tfrac{11\pi}{6}$ 
				\item $t=\tfrac{5\pi}{3}$ 
			\end{enumerate}
		\end{multicols}\vspace{-2.75em}
	\end{spacing}	
\end{exercise}
\begin{example}[modèle] Pour chaque $t$, retrouver $-\pi< t'\leqslant \pi$ tel que $t=t'+2k\pi$ ou $k\in\mathbb{Z}$ et déduire les coordonnées du point associé à $t$. 
\end{example}
\begin{proof}[solution]Lire la solution. Pouquoi avoir entouré $4\pi$ et $-10\pi$ ?\hfill\\  
\parbox{0.45\linewidth}{
$\begin{WithArrows}
	t& =\frac{25\pi}{6}\\ 
	3\pi<\frac{25\pi}{6}&\leqslant \tikz[baseline =(A.base),remember picture]\node[draw,ellipse, color=red, line width=2pt](A) {$4\pi$};\Arrow{$-4\pi$}\\
	-\pi <\frac{25\pi}{6}-4\pi &\leqslant 0\\
	-\pi <\frac{\pi}{6} &\leqslant 0
\end{WithArrows}$ 
}
\hfill 
\parbox{0.45\linewidth}{
$\begin{WithArrows}
	t& =-\frac{39\pi}{4}\\ 
	\tikz[baseline =(A.base),remember picture]\node[draw,ellipse, color=red, line width=2pt](A) {$-10\pi$};<-\frac{39\pi}{4}&\leqslant -9\pi\Arrow{$+10\pi$}\\
	0 <-\frac{39\pi}{4}+10\pi &\leqslant \pi\\
	0 < \frac{\pi}{4} &\leqslant \pi
\end{WithArrows}$  
} \\
Le point associé à $\frac{25\pi}{6}$ est $P(\tfrac{\sqrt{3}}{2}~;~\tfrac{1}{2})$.
	\hfill 
	Le point associé à $\frac{\pi}{4}$ est $P(\tfrac{\sqrt{2}}{2}~;~\tfrac{\sqrt{2}}{2})$.	
\end{proof}
\begin{exercise}[à vous]  Mêmes consignes
	\vspace{-0.75em}
	\begin{spacing}{1.8}
		\begin{multicols}{5} 
			\begin{enumerate}[label={\arabic*)},leftmargin=*]
				\item $t=\tfrac{13\pi}{6}$ 
				\item $t=\tfrac{41\pi}{6}$ 
				\item $t=-\tfrac{11\pi}{3}$ 
				\item $t=\tfrac{31\pi}{6}$ 
				\item $t=-\tfrac{41\pi}{4}$ 
			\end{enumerate}
		\end{multicols}\vspace{-2.75em}
	\end{spacing}	
\end{exercise}


 

 




%----------------------------------------------------------------------------------------
%	 
%----------------------------------------------------------------------------------------
\newpage 
\loadgeometry{marginlayout}  
\pagestyle{scrheadings} % Print headers again  
\KOMAoptions{
	headwidth=textwithmarginpar,
	headsepline=true,
	headsepline= 0.5pt,
	footwidth=textwithmarginpar,
}  
\onehalfspacing 
\section{Mesure d'angles orientés}  
\begin{marginfigure} 
	\centering\noindent
	\begin{tikzpicture}[scale=1.5];
		\tkzSetUpPoint[size=3, fill=black,] %shape=cross out
		\tkzfctset{tan style/.style={->,>=latex, color=red, line width=1.2pt}}   
		\tkzSetUpLine[line width= 1.2pt, add=.2 and .2] 
		%		
		\tkzDefPoint(0,0){O}  
		\tkzDefPoint(1,0){I}
		\tkzDefShiftPoint[O](-15:1){A'}
		\tkzDefShiftPoint[O](-15:1.3){A}
		\tkzDefShiftPoint[O](55:1){B'}
		\tkzDefShiftPoint[O](55:1.7){B}
		\tkzFillAngle[size=1.3,left color=gray!70,
		right color=white](A,O,B)

		\tkzDrawArc[angles,color=blue,dashed,line width=1.2pt](O,I)(55,-15)
		\tkzDrawArc[angles,color=red, line width=1.2pt, -,>=latex](O,I)(-15,55)
		\tkzDrawSegments[line width=1pt,-,>=latex](O,A O,B)
		\tkzMarkAngle[size=0.4](A,O,B)
		\tkzLabelAngle[pos=0.5](A,O,B){$\theta$}
		\tkzLabelSegment[  below, font=\footnotesize](O,A'){$1$}
		\tkzDrawPoints[](A',B')
		\tkzLabelPoints[right](A,B)
		\tkzLabelPoints[left](O)
		\tkzText[](1.2,0.5){$L$}
	\end{tikzpicture} 
	\caption{L'angle $\widehat{AOB}$ est une partie du plan délimité par deux deux demi-droites $[OA)$ et $[OB)$.}
\end{marginfigure} 
\begin{definition}[le radian] On trace un cercle de rayon $1$.  
	La mesure en~\si{\rad} de l'angle $\widehat{AOB}$ est égale à longueur de l'arc intercepté. 
\end{definition}  
\noindent\parbox{1\linewidth}{\centering
\null\hfill
\begin{tikzpicture}[scale=1.2];
	\tkzSetUpPoint[size=3, fill=black,] %shape=cross out
	\tkzfctset{tan style/.style={->,>=latex, color=red, line width=1.2pt}}   
	\tkzSetUpLine[line width= 1.2pt, add=.2 and .2] 
	%		
	\tkzDefPoint(0,0){O}  
	\tkzDefPoint(1,0){I}
	\tkzDefShiftPoint[O](-15:1){A}
	\tkzDefShiftPoint[O](42:1){B}
	\tkzFillAngle[size=1.0,left color=gray!70,
	right color=white](A,O,B)

	\tkzDrawArc[angles,color=blue,dashed,line width=1.2pt](O,I)(42,-15)
	\tkzDrawArc[angles,color=red, line width=1.2pt, -,>=latex](O,I)(-15,42)
	\tkzDrawSegments[line width=1pt,-,>=latex](O,A O,B)
	\tkzMarkAngle[size=0.4](A,O,B)
	\tkzLabelAngle[pos=0.8,font=\scriptsize](A,O,B){$\SI{1}{\rad}$}
	\tkzLabelSegment[  below, font=\footnotesize](O,A'){$1$}
	%		\tkzLabelPoints[right](A,B)
	\tkzLabelPoints[left](O)
	\tkzText[](1.2,0.4){$1$}
\end{tikzpicture}  
\hfill 
\begin{tikzpicture}[scale=1.2];
	\tkzSetUpPoint[size=3, fill=black,] %shape=cross out
	\tkzfctset{tan style/.style={->,>=latex, color=red, line width=1.2pt}}   
	\tkzSetUpLine[line width= 1.2pt, add=.2 and .2] 
	%		
	\tkzDefPoint(0,0){O}  
	\tkzDefPoint(1,0){I}
	\tkzDefShiftPoint[O](-15:1){A}
	\tkzDefShiftPoint[O](99.6:1){B}
	\tkzFillAngle[size=1.0,left color=gray!70,
	right color=white](A,O,B)

	\tkzDrawArc[angles,color=blue,dashed,line width=1.2pt](O,I)(99.6,-15)
	\tkzDrawArc[angles,color=red, line width=1.2pt, -,>=latex](O,I)(-15,99.6)
	\tkzDrawSegments[line width=1pt,-,>=latex](O,A O,B)
	\tkzMarkAngle[size=0.4](A,O,B)
	\tkzLabelAngle[pos=0.6,font=\scriptsize](A,O,B){$\SI{2}{\rad}$}
	\tkzLabelSegment[  below, font=\footnotesize](O,A'){$1$}
	%		\tkzLabelPoints[right](A,B)
	\tkzLabelPoints[left](O)
	\tkzText[](1.0,0.6){$2$}
\end{tikzpicture}  
\hfill
\begin{tikzpicture}[scale=1.2];
	\tkzSetUpPoint[size=3, fill=black,] %shape=cross out
	\tkzfctset{tan style/.style={->,>=latex, color=red, line width=1.2pt}}   
	\tkzSetUpLine[line width= 1.2pt, add=.2 and .2] 
	%		
	\tkzDefPoint(0,0){O}  
	\tkzDefPoint(1,0){I}
	\tkzDefShiftPoint[O](-15:1){A}
	\tkzDefShiftPoint[O](157:1){B}
	\tkzFillAngle[size=1.0,left color=gray!70,
	right color=white](A,O,B)

	\tkzDrawArc[angles,color=blue,dashed,line width=1.2pt](O,I)(157,-15)
	\tkzDrawArc[angles,color=red, line width=1.2pt, -,>=latex](O,I)(-15,157)
	\tkzDrawSegments[line width=1pt,-,>=latex](O,A O,B)
	\tkzMarkAngle[size=0.4](A,O,B)
	\tkzLabelAngle[pos=0.6,font=\scriptsize](A,O,B){$\SI{3}{\rad}$}
	\tkzLabelSegment[  below, font=\footnotesize](O,A'){$1$}
	%		\tkzLabelPoints[right](A,B)
	\tkzLabelPoints[below](O)
	\tkzText[](1.0,0.6){$3$}
\end{tikzpicture} 
\hfill
\null
}
\begin{eBox}
	Un angle plat correspond à \ang{180}, soit une mesure en radian  de $\pi$ :
	\setlength{\abovedisplayskip}{3pt}
	\setlength{\belowdisplayskip}{3pt} 
	$$
	\ang{180}=\pi~\si{\rad}
	%\qquad \SI{1}{\rad}=\left(\frac{180}{\pi}\right)^\circ
	%\qquad \ang{1}=\frac{\pi}{180}~\si{\rad}
	$$
\end{eBox}
\begin{example}[convertir~\si{\rad} $\leftrightarrow$ degrés]\label{chapter08:example:conversion}  $\ang{60}$ en~\si{\rad}, et $\tfrac{\pi}{6}~\si{\rad}$ en degrés. \\
	\null\hfill 
	\renewcommand{\arraystretch}{1.25}
	\begin{tabularx}{0.8\linewidth}{|>{\columncolor{blue!20}}c|*{6}{>{\centering \arraybackslash}X|}}
		\hline
		mesure en degrés$^\circ$ & $360$ & $180$ & $90$ & $60$ & $45$ & \\\hline
		{mesure en~\si{\rad}} & $2\pi$ & $\pi$ & $\tfrac{\pi}{2}$ & &  & $\tfrac{\pi}{6}$ \\\hline
	\end{tabularx} 
	\hfill\null
\end{example} 
\begin{proof}[solution]  
	$\ang{60}=60\times \tfrac{\pi}{180}~\si{\rad}=\tfrac{\pi}{3}~\si{\rad}$ et $\tfrac{\pi}{6}~\si{\rad}=\tfrac{\pi}{6} \times \frac{180}{\pi}=\ang{30}$. 
\end{proof} 
\begin{remark}
	Pour référence : $\SI{1}{\rad}\approx \ang{57.296}$ et $\ang{1}\approx\numprint{0.01745}~\si{\rad}$
\end{remark}
\begin{property}[longueur d'un arc] 
\begin{marginfigure}
		\centering
		\begin{tikzpicture}[scale=1.8];
			\tkzSetUpPoint[size=3, fill=black,] %shape=cross out
			\tkzfctset{tan style/.style={->,>=latex, color=red, line width=1.2pt}}   
			\tkzSetUpLine[line width= 1.2pt, add=.2 and .2] 
			%		
			\tkzDefPoint(0,0){O}  
			\tkzDefPoint(1,0){I}
			\tkzDefShiftPoint[O](-15:1){A}
			\tkzDefShiftPoint[O](55:1){B}
			\tkzFillAngle[size=1.0,left color=gray!70,
			right color=white](A,O,B)

			\tkzDrawArc[angles,color=blue,line width=1.2pt](O,I)(55,-15)
			\tkzDrawArc[angles,color=red, line width=1.2pt, -,>=latex](O,I)(-15,55)
			\tkzDrawSegments[line width=1pt,-,>=latex](O,A O,B)
			\tkzMarkAngle[size=0.4](A,O,B)
			\tkzLabelAngle[pos=0.5](A,O,B){$\theta$}
			\tkzLabelSegment[  below, font=\footnotesize](O,A'){$r$}
			%		\tkzLabelPoints[right](A,B)
			\tkzLabelPoints[left](O)
			\tkzText[text width=2.5cm](1.8,0.5){$s$}
		\end{tikzpicture}
		\caption{$s=r\theta$.}
	\end{marginfigure}  
Pour un cercle de rayon $r$. La longueur $s$ d'un arc d'\textbf{angle au centre} $\theta$ (exprimé en~\si{\rad}) est :
	\setlength{\abovedisplayskip}{3pt}
	\setlength{\belowdisplayskip}{3pt} 
	$$
	s=r\theta
	$$  
\end{property}
\begin{proof} $s=\tfrac{\theta}{2\pi}\times \text{périmètre}=\tfrac{\theta}{2\pi}\times 2\pi r=r\theta$.
\end{proof}
\begin{example}[longueur d'un arc $\leftrightarrow$ angle au centre] \hfill 
\begin{enumerate}[label=\alph*),leftmargin=*]
	\item Détermine la longueur de l'arc d'un cercle de rayon \SI{10}{\meter} d'angle au centre \ang{30}
	\item Détermine l'angle au centre en~\si{\rad} et en degrés pour un arc   de longueur \SI{6}{\meter} sur un cercle de rayon \SI{4}{\meter}.
\end{enumerate} 
\end{example}
\begin{proof}[solution]  a) D'après \ref{chapter08:example:conversion} :  $\ang{30}=\tfrac{\pi}{6}~\si{\rad}$ d'où $s=r\theta=10\times \tfrac{\pi}{6}= \tfrac{5\pi}{3}~\si{\meter}$\\
	b) $\theta=\tfrac{s}{r}=\tfrac{6}{4}=\frac{3}{2}~\si{\rad}=\tfrac{3}{2}\times \tfrac{180}{\pi}\approx \ang{86} $.
\end{proof}
\begin{definition}
\begin{marginfigure}%[!h] 
	\centering
\begin{tikzpicture}[scale=0.75];
\tkzSetUpPoint[size=3, fill=black,shape=cross out]
\tkzfctset{tan style/.style={->,>=latex, color=red, line width=1.2pt}}   
\tkzSetUpLine[line width= 1.2pt, add=.2 and .2] 

%\tkzInit[xmin=-4.0,xmax=4.0,ymin=-4.0,ymax=4.0,xstep=1.0, ystep=1.0];  
	\tkzDefPoint(0,0){O} 
	\tkzDefPoint(2,1){A}
	\tkzDefPoint(-2,1){B}
	\tkzLabelPoints[above](A, B)
	\tkzLabelPoints[below,font=\small](O)
	\tkzInterLC[R](O,A)(O,1cm) \tkzGetPoints{M2}{M1}
	\tkzDrawSegments[->,>=latex](O,A O,B)
	\tkzMarkAngle[->,>=latex,size=0.5,line width=1.2pt](A,O,B) 
%	\draw[-latex, line width = 2] (M1) arc (26:155:0.5) ;
\begin{scope}[yshift=2cm]
	\tkzDefPoint(0,0){O} 
	\tkzDefPoint(2,1){A}
	\tkzDefPoint(-2,1){B}
	\tkzLabelPoints[above](A)
	\tkzLabelPoints[above](B)
	\tkzLabelPoints[above,font=\small](O)
	\tkzInterLC[R](O,A)(O,1cm) \tkzGetPoints{M2}{M1}
	\tkzDrawSegments[->,>=latex](O,A O,B)
	\tkzMarkAngle[<-,>=latex,size=0.5,line width=1.2pt](B,O,A) 
\end{scope}
\end{tikzpicture}
\caption{Angle orienté $(\protect\vv{OA}~;~\protect\vv{OB})$ avec une mesure positive ou négative}
\end{marginfigure}
 L'angle orienté $(\vv{OR_1}~;~\vv{OR_2})$ est défini par un sommet $O$, une demi-droite $[OR_1)$ initiale, et une demi-droite finale $[OR_2)$. 
\end{definition}  
Pour mesurer un angle orienté, on regarde l'angle de la rotation qui transforme $[OR_1)$ en $[OR_2)$. Si la rotation se fait dans le sens positif, la mesure  de l'angle est positive.\\
Si la rotation  se fait dans le sens négatif, la mesure de l'angle est négative. 
\begin{figure}[!h]
	\begin{tikzpicture}[scale=1.0];
		\tkzSetUpPoint[size=3, fill=black,] %shape=cross out
		\tkzfctset{tan style/.style={->,>=latex, color=red, line width=1.2pt}}   
		\tkzSetUpLine[line width= 1.2pt, add=.2 and .2] 
		%		
		\tkzDefPoint(0,0){O}  
		\tkzDefPoint(1,0){I}
		\tkzDefShiftPoint[O](30:1){A'}
		\tkzDefShiftPoint[O](30:1.5){A}
		\tkzDefShiftPoint[O](65:1){B'}
		\tkzDefShiftPoint[O](65:1.9){B}
		\tkzDefShiftPoint[O](50:1.00){T}
		\tkzLabelPoint[above right](T){\scriptsize\ang{30}}
%		\tkzDrawArc[angles,color=blue,dashed,line width=1.2pt](O,I)(65,35)
		\tkzDrawArc[angles,color=red, line width=1.2pt, ->,>=latex](O,I)(30,65)
		\tkzDrawSegments[line width=1pt,->,>=latex](O,A O,B)
%		\tkzMarkAngle[size=0.2](A,O,B)
%		\tkzLabelAngle[pos=0.4](A,O,B){$\theta$}
%		\tkzLabelSegment[  left, font=\footnotesize](O,A'){$1$}
%		\tkzDrawPoints[](A',B')
		\tkzLabelPoints[above right](A,B)
		\tkzLabelPoints[below left](O) 
	\end{tikzpicture}
	\begin{tikzpicture}[scale=1.0];
	\tkzSetUpPoint[size=3, fill=black,] %shape=cross out
	\tkzfctset{tan style/.style={->,>=latex, color=red, line width=1.2pt}}   
	\tkzSetUpLine[line width= 1.2pt, add=.2 and .2] 
	%
	\tkzInit[xmin=-1.1,xmax=1.1,ymin=-1.0,ymax=1.0,xstep=1.0, ystep=1.0];  
		
	\tkzDefPoint(0,0){O}  
	\tkzDefPoint(1,0){I}
	\tkzDefShiftPoint[O](30:1){A'}
	\tkzDefShiftPoint[O](30:1.5){A}
	\tkzDefShiftPoint[O](65:1){B'}
	\tkzDefShiftPoint[O](65:1.9){B}
	\tkzDefShiftPoint[O](-50:1.00){T}
	\tkzLabelPoint[above right](T){\scriptsize\ang{390}}
	%		\tkzDrawArc[angles,color=blue,dashed,line width=1.2pt](O,I)(65,35)
%	\tkzDrawArc[angles,color=red, line width=1.2pt, ->,>=latex](O,I)(30,65)
	\tkzDrawSegments[line width=1pt,->,>=latex](O,A O,B)
	%		\tkzMarkAngle[size=0.2](A,O,B)
	%		\tkzLabelAngle[pos=0.4](A,O,B){$\theta$}
	%		\tkzLabelSegment[  left, font=\footnotesize](O,A'){$1$}
	%		\tkzDrawPoints[](A',B')
	\tkzLabelPoints[above right](A,B)
	\tkzLabelPoints[below left](O) 
\tkzFctPar[smooth,samples=50,domain=0.5236:7.414,->,>=latex,line width=0.8pt,color=red]{(0.35+6*t/74.14)*cos(t)}{(0.35+6*t/74.14)*sin(t)} 
\end{tikzpicture} 
\begin{tikzpicture}[scale=1.0];
\tkzSetUpPoint[size=3, fill=black,] %shape=cross out
\tkzfctset{tan style/.style={->,>=latex, color=red, line width=1.2pt}}   
\tkzSetUpLine[line width= 1.2pt, add=.2 and .2] 
%
\tkzInit[xmin=-1.1,xmax=1.1,ymin=-1.0,ymax=1.0,xstep=1.0, ystep=1.0];  

\tkzDefPoint(0,0){O}  
\tkzDefPoint(1,0){I}
\tkzDefShiftPoint[O](30:1){A'}
\tkzDefShiftPoint[O](30:1.5){A}
\tkzDefShiftPoint[O](65:1){B'}
\tkzDefShiftPoint[O](65:1.9){B}
\tkzDefShiftPoint[O](-50:1.00){T}
\tkzLabelPoint[above right](T){\scriptsize$-\ang{330}$}
%		\tkzDrawArc[angles,color=blue,dashed,line width=1.2pt](O,I)(65,35)
%	\tkzDrawArc[angles,color=red, line width=1.2pt, ->,>=latex](O,I)(30,65)
\tkzDrawSegments[line width=1pt,->,>=latex](O,A O,B)
%		\tkzMarkAngle[size=0.2](A,O,B)
%		\tkzLabelAngle[pos=0.4](A,O,B){$\theta$}
%		\tkzLabelSegment[  left, font=\footnotesize](O,A'){$1$}
%		\tkzDrawPoints[](A',B')
\tkzLabelPoints[above right](A,B)
\tkzLabelPoints[below left](O) 
\tkzFctPar[smooth,samples=50,domain=-0.5236:5.148,->,>=latex,line width=0.8pt,color=red]{(0.5+2*t/52)*cos(-t)}{(0.5+2*t/52)*sin(-t)} 
\end{tikzpicture}
	\caption{L'angle orienté $(\protect\vv{OA}~;~\protect\vv{OB})$ a plusieurs mesures.}
\end{figure}


\begin{marginfigure}
\begin{tikzpicture}[scale=1.2];
	\tkzSetUpPoint[size=3, fill=black,] %shape=cross out
	\tkzfctset{tan style/.style={->,>=latex, color=red, line width=1.2pt}}   
	\tkzSetUpLine[line width= 1.2pt, add=.2 and .2] 
	%		
	\tkzDefPoint(0,0){O}  
	\tkzDefPoint(1,0){I}
	\tkzDefShiftPoint[O](-45:1){A'}
	\tkzDefShiftPoint[O](-45:1.6){A}
	\tkzDefShiftPoint[O](65:1){B'}
	\tkzDefShiftPoint[O](65:1.95){B}
	\tkzDrawArc[angles,color=blue,dashed,line width=1.2pt](O,I)(65,-45)
	\tkzDrawArc[angles,color=red, line width=1.2pt, ->,>=latex](O,I)(-45,65)
	\tkzDrawSegments[line width=1pt,->,>=latex](O,A O,B)
	\tkzMarkAngle[size=0.2](A,O,B)
	\tkzLabelAngle[pos=0.4](A,O,B){$\theta$}
	\tkzLabelSegment[  left, font=\footnotesize](O,A'){$1$}
	\tkzDrawPoints[](A',B')
	\tkzLabelPoints[right](A,B)
	\tkzLabelPoints[left](O)
	\tkzText[text width=2.5cm](2.2,0){Mesure en radian de $\theta$}
\end{tikzpicture}
\caption{}
\end{marginfigure}
\begin{definition}[mesure principale en radian d'un angle  $(\vv{OA}~;~\vv{OB})$] On trace le cercle de centre $O$ et de rayon $1$. \\
La \textbf{mesure principale} $\theta$ de l'angle  $(\vv{OA}~;~\vv{OB})$ est la longueur du plus court arc  de cercle orienté intercepté. \\
La \textbf{mesure principale} vérifie $-\pi < \theta\leqslant \pi$. \\
L'angle $(\vv{OA}~;~\vv{OB})$ a une infinité de mesures de la forme $\theta + 2k\pi$ ($k\in\mathbb{Z}$). Le $k$ détermine en fait un nombre de tours que l’on effectue sans le sens direct si $k>0$, et dans le sens indirect si $k<0$.
\end{definition}  
\begin{figure*}[!ht]
\begin{tikzpicture}[scale=1.25]
	\tkzSetUpPoint[size=2, fill=black]%,shape=cross out
	\tkzfctset{tan style/.style={->,>=latex, color=red, line width=1.2pt}}   
	\tkzSetUpLine[line width= 1.2pt, add=.2 and .2]  
	\tkzInit[xmin=-1.0,xmax=1.0,ymin=-1.0,ymax=1.0,xstep=1.0, ystep=1.0];  
	
	%\tkzDrawX[ font=\small, line width=2pt, label=$x$ , below=5pt ,left space=0.6,right space=0.6] 
	%\tkzDrawY[ font=\small, line width=2pt , label={$y$} , right =5pt, down space=0.5,up space=0.5]  
	%\tkzLabelXY[orig = false,  font=\scriptsize ] 
	%\tkzRep[color  = Maroon,xlabel=,ylabel=, font=\scriptsize, line width = 2pt,xnorm=1, ynorm=1]   
	\def\a{1} 
	\tkzDefPoint(0,0){O}  
	\tkzDefPoint(1,0){I}  
	\tkzLabelPoint[below left](O){$0$}
	%	\tkzLabelPoint[below right](I){$1$}
	\tkzLabelSegment[below](O,I){$1$}
	\tkzDrawCircle[color=blue, line width=1.2pt](O,I)
	\tkzDrawArc[color=red,->,>=latex,angles,line width=2pt](O,I)(0,{\a*180/3.14159})
	\tkzDefShiftPoint[O]({\a*180/3.14159}:1){P}
	\tkzDrawPoints(O,P) 
	\tkzDefShiftPoint[O]({\a*90/3.14159}:1){T}	
	\tkzLabelPoint[above right](T){\small\color{red}$\SI{1}{\rad}$}	 
	\tkzDrawLine[line width=1.2pt,add=.0 and .4,->,>=latex,color=red](O,P)
	\tkzDrawLine[line width=1.2pt,add=.0 and .4,->,>=latex](O,I)
\end{tikzpicture} 
\begin{tikzpicture}[scale=1.25]
	\tkzSetUpPoint[size=2, fill=black]%,shape=cross out
	\tkzfctset{tan style/.style={->,>=latex, color=red, line width=1.2pt}}   
	\tkzSetUpLine[line width= 1.2pt, add=.2 and .2]  
	\tkzInit[xmin=-1.0,xmax=1.0,ymin=-1.0,ymax=1.0,xstep=1.0, ystep=1.0];  
	
	%\tkzDrawX[ font=\small, line width=2pt, label=$x$ , below=5pt ,left space=0.6,right space=0.6] 
	%\tkzDrawY[ font=\small, line width=2pt , label={$y$} , right =5pt, down space=0.5,up space=0.5]  
	%\tkzLabelXY[orig = false,  font=\scriptsize ] 
	%\tkzRep[color  = Maroon,xlabel=,ylabel=, font=\scriptsize, line width = 2pt,xnorm=1, ynorm=1]   
	\def\a{2} 
	\tkzDefPoint(0,0){O}  
	\tkzDefPoint(1,0){I}  
	\tkzLabelPoint[below left](O){$0$}
	%	\tkzLabelPoint[below right](I){$1$}
	\tkzLabelSegment[below](O,I){$1$}
	\tkzDrawCircle[color=blue, line width=1.2pt](O,I)
	\tkzDrawArc[color=red,->,>=latex,angles,line width=2pt](O,I)(0,{\a*180/3.14159})
	\tkzDefShiftPoint[O]({\a*180/3.14159}:1){P}
	\tkzDrawPoints(O, P) 
	\tkzDefShiftPoint[O]({\a*90/3.14159}:1){T}	
	\tkzLabelPoint[above right](T){\small\color{red}$\SI{2}{\rad}$}	 
	\tkzDrawLine[line width=1.2pt,add=.0 and .4,->,>=latex,color=red](O,P)
	\tkzDrawLine[line width=1.2pt,add=.0 and .4,->,>=latex](O,I)
\end{tikzpicture} 
\begin{tikzpicture}[scale=1.25]
	\tkzSetUpPoint[size=2, fill=black]%,shape=cross out
	\tkzfctset{tan style/.style={->,>=latex, color=red, line width=1.2pt}}   
	\tkzSetUpLine[line width= 1.2pt, add=.2 and .2]  
	\tkzInit[xmin=-1.0,xmax=1.0,ymin=-1.0,ymax=1.0,xstep=1.0, ystep=1.0];  
	
	%\tkzDrawX[ font=\small, line width=2pt, label=$x$ , below=5pt ,left space=0.6,right space=0.6] 
	%\tkzDrawY[ font=\small, line width=2pt , label={$y$} , right =5pt, down space=0.5,up space=0.5]  
	%\tkzLabelXY[orig = false,  font=\scriptsize ] 
	%\tkzRep[color  = Maroon,xlabel=,ylabel=, font=\scriptsize, line width = 2pt,xnorm=1, ynorm=1]   
	\def\a{3} 
	\tkzDefPoint(0,0){O}  
	\tkzDefPoint(1,0){I}  
	\tkzLabelPoint[below left](O){$0$}
	%	\tkzLabelPoint[below right](I){$1$}
	\tkzLabelSegment[below](O,I){$1$}
	\tkzDrawCircle[color=blue, line width=1.2pt](O,I)
	\tkzDrawArc[color=red,->,>=latex,angles,line width=2pt](O,I)(0,{\a*180/3.14159})
	\tkzDefShiftPoint[O]({\a*180/3.14159}:1){P}
	\tkzDrawPoints(O, P) 
	\tkzDefShiftPoint[O]({\a*90/3.14159}:1){T}	
	\tkzLabelPoint[above right](T){\small\color{red}$\SI{3}{\rad}$}	 
	\tkzDrawLine[line width=1.2pt,add=.0 and .4,->,>=latex,color=red](O,P)
	\tkzDrawLine[line width=1.2pt,add=.0 and .4,->,>=latex](O,I)
\end{tikzpicture} 
\begin{tikzpicture}[scale=1.25]
	\tkzSetUpPoint[size=2, fill=black]%,shape=cross out
	\tkzfctset{tan style/.style={->,>=latex, color=red, line width=1.2pt}}   
	\tkzSetUpLine[line width= 1.2pt, add=.2 and .2]  
	\tkzInit[xmin=-1.0,xmax=1.0,ymin=-1.0,ymax=1.0,xstep=1.0, ystep=1.0];  
	
	%\tkzDrawX[ font=\small, line width=2pt, label=$x$ , below=5pt ,left space=0.6,right space=0.6] 
	%\tkzDrawY[ font=\small, line width=2pt , label={$y$} , right =5pt, down space=0.5,up space=0.5]  
	%\tkzLabelXY[orig = false,  font=\scriptsize ] 
	%\tkzRep[color  = Maroon,xlabel=,ylabel=, font=\scriptsize, line width = 2pt,xnorm=1, ynorm=1]   
	\def\a{3.14159} 
	\tkzDefPoint(0,0){O}  
	\tkzDefPoint(1,0){I}  
	\tkzLabelPoint[below left](O){$0$}
	%	\tkzLabelPoint[below right](I){$1$}
	\tkzLabelSegment[below](O,I){$1$}
	\tkzDrawCircle[color=blue, line width=1.2pt](O,I)
	\tkzDrawArc[color=red,->,>=latex,angles,line width=2pt](O,I)(0,{\a*180/3.14159})
	\tkzDefShiftPoint[O]({\a*180/3.14159}:1){P}
	\tkzDrawPoints(O,P) 
	\tkzDefShiftPoint[O]({\a*90/3.14159}:1){T}	
	\tkzLabelPoint[above right](T){\small\color{red}$\pi~\si{\rad}$}	 
	\tkzDrawLine[line width=1.2pt,add=.0 and .4,->,>=latex,color=red](O,P)
	\tkzDrawLine[line width=1.2pt,add=.0 and .4,->,>=latex](O,I)
\end{tikzpicture} 
\begin{tikzpicture}[scale=1.25]
	\tkzSetUpPoint[size=2, fill=black]%,shape=cross out
	\tkzfctset{tan style/.style={->,>=latex, color=red, line width=1.2pt}}   
	\tkzSetUpLine[line width= 1.2pt, add=.2 and .2]  
	\tkzInit[xmin=-1.0,xmax=1.0,ymin=-1.0,ymax=1.0,xstep=1.0, ystep=1.0];  
	
	%\tkzDrawX[ font=\small, line width=2pt, label=$x$ , below=5pt ,left space=0.6,right space=0.6] 
	%\tkzDrawY[ font=\small, line width=2pt , label={$y$} , right =5pt, down space=0.5,up space=0.5]  
	%\tkzLabelXY[orig = false,  font=\scriptsize ] 
	%\tkzRep[color  = Maroon,xlabel=,ylabel=, font=\scriptsize, line width = 2pt,xnorm=1, ynorm=1]   
	\def\a{3.14159/2} 
	\tkzDefPoint(0,0){O}  
	\tkzDefPoint(1,0){I}  
	\tkzLabelPoint[below left](O){$0$}
	%	\tkzLabelPoint[below right](I){$1$}
	\tkzLabelSegment[below](O,I){$1$}
	\tkzDrawCircle[color=blue, line width=1.2pt](O,I)
	\tkzDrawArc[color=red,->,>=latex,angles,line width=2pt](O,I)(0,{\a*180/3.14159})
	\tkzDefShiftPoint[O]({\a*180/3.14159}:1){P}
	\tkzDrawPoints(O,P) 
	\tkzDefShiftPoint[O]({\a*90/3.14159}:1){T}	
	\tkzLabelPoint[above right](T){\small\color{red}$\tfrac{\pi}{2}~\si{\rad}$}	 
	\tkzDrawLine[line width=1.2pt,add=.0 and .4,->,>=latex,color=red](O,P)
	\tkzDrawLine[line width=1.2pt,add=.0 and .4,->,>=latex](O,I)
\end{tikzpicture} 
\caption{Dans un plan repéré et orienté, il est usuel de représenter un angle orienté en prenant comme demi-droite initiale l'axe $[Ox)$ des abscisses.}
\end{figure*}
\begin{figure}[!ht]
	\begin{tikzpicture}[scale=1.1]
	\tkzSetUpPoint[size=2, fill=black]%,shape=cross out
	\tkzfctset{tan style/.style={->,>=latex, color=red, line width=1.2pt}}   
	\tkzSetUpLine[line width= 1.2pt, add=.2 and .2]  
	\tkzInit[xmin=-1.1,xmax=1.1,ymin=-1.0,ymax=1.0,xstep=1.0, ystep=1.0];  
	
\tkzDrawX[ noticks,font=\small, line width=2pt, label=$x$ , below=5pt ,left space=0.4,right space=0.5] 
	\tkzDrawY[ noticks,font=\small, line width=2pt , label={$y$} , right =5pt, down space=0.6,up space=0.5]  
	%\tkzLabelXY[orig = false,  font=\scriptsize ] 
	%\tkzRep[color  = Maroon,xlabel=,ylabel=, font=\scriptsize, line width = 2pt,xnorm=1, ynorm=1]   
	\def\a{-5*3.14159/6} 
	\tkzDefPoint(0,0){O}  
	\tkzDefPoint(1,0){I}  
	\tkzLabelPoint[below left](O){$0$}
	\tkzLabelPoint[below right](I){$A$}
	%		\tkzLabelSegment[below](O,I){$1$}
	\tkzDrawCircle[color=blue, line width=1.2pt](O,I)
	\tkzDrawArc[color=red,<-,>=latex,angles,line width=2pt](O,I)({\a*180/3.14159},0) 
	\tkzDefShiftPoint[O]({\a*180/3.14159}:1){P}
	\tkzDrawPoints(O,P) 
	\tkzDefShiftPoint[O]({\a*180/3.14159}:1.6){B}
	\tkzLabelPoint[below](B){$B$}
	\tkzDefShiftPoint[O]({\a*90/3.14159}:1){T}	
	\tkzLabelPoint[below](T){\scriptsize\color{red}mesure principale $\tfrac{-5\pi}{6}~\si{\rad}$}	 
	\tkzDrawLine[line width=1.2pt,add=.0 and .6,->,>=latex,color=red](O,P)
	%		\tkzDrawLine[line width=1.2pt,add=.0 and .4,->,>=latex](O,I)
\end{tikzpicture} 
	\begin{tikzpicture}[scale=1.1]
	\tkzSetUpPoint[size=2, fill=black]%,shape=cross out
	\tkzfctset{tan style/.style={->,>=latex, color=red, line width=1.2pt}}   
	\tkzSetUpLine[line width= 1.2pt, add=.2 and .2]  
	\tkzInit[xmin=-1.1,xmax=1.1,ymin=-1.0,ymax=1.0,xstep=1.0, ystep=1.0]; 
		
\tkzDrawX[ noticks,font=\small, line width=2pt, label=$x$ , below=5pt ,left space=0.4,right space=0.5] 
	\tkzDrawY[ noticks,font=\small, line width=2pt , label={$y$} , right =5pt, down space=0.6,up space=0.5]  
		%\tkzLabelXY[orig = false,  font=\scriptsize ] 
		%\tkzRep[color  = Maroon,xlabel=,ylabel=, font=\scriptsize, line width = 2pt,xnorm=1, ynorm=1]   
		\def\a{7*3.14159/6} 
		\tkzDefPoint(0,0){O}  
		\tkzDefPoint(1,0){I}  
		\tkzLabelPoint[above left](O){$0$}
		\tkzLabelPoint[below right](I){$A$}
%		\tkzLabelSegment[below](O,I){$1$}
		\tkzDrawCircle[color=blue, line width=1.2pt](O,I)
		\tkzDrawArc[color=red,->,>=latex,angles,line width=2pt](O,I)(0,{\a*180/3.14159}) 
		\tkzDefShiftPoint[O]({\a*180/3.14159}:1){P}
		\tkzDrawPoints(O,P)  
		\tkzDefShiftPoint[O]({\a*180/3.14159}:1.6){B}
		\tkzLabelPoint[below](B){$B$}
		\tkzDefShiftPoint[O]({\a*90/3.14159}:1){T}	
		\tkzLabelPoint[above  ](T){\scriptsize\color{red}$\tfrac{7\pi}{6}=-\tfrac{5\pi}{6}+2\pi$}	 
		\tkzDrawLine[line width=1.2pt,add=.0 and .6,->,>=latex,color=red](O,P)
%		\tkzDrawLine[line width=1.2pt,add=.0 and .4,->,>=latex](O,I)
	\end{tikzpicture} 
	\begin{tikzpicture}[scale=1.1]
	\tkzSetUpPoint[size=2, fill=black]%,shape=cross out
	\tkzfctset{tan style/.style={->,>=latex, color=red, line width=1.2pt}}   
	\tkzSetUpLine[line width= 1.2pt, add=.2 and .2]  
	\tkzInit[xmin=-1.1,xmax=1.1,ymin=-1.0,ymax=1.0,xstep=1.0, ystep=1.0]; 

\tkzDrawX[ noticks,font=\small, line width=2pt, label=$x$ , below=5pt ,left space=0.4,right space=0.5] 
	\tkzDrawY[ noticks,font=\small, line width=2pt , label={$y$} , right =5pt, down space=0.6,up space=0.5]  
%\tkzLabelXY[orig = false,  font=\scriptsize ] 
%\tkzRep[color  = Maroon,xlabel=,ylabel=, font=\scriptsize, line width = 2pt,xnorm=1, ynorm=1]   
\def\a{7*3.14159/6} 
\tkzDefPoint(0,0){O}  
\tkzDefPoint(1,0){I}  
\tkzLabelPoint[below left](O){$0$}
\tkzLabelPoint[below right](I){$A$}
%		\tkzLabelSegment[below](O,I){$1$}
\tkzDrawCircle[color=blue, line width=1.2pt](O,I)
%\tkzDrawArc[color=red,->,>=latex,angles,line width=2pt](O,I)(0,{\a*180/3.14159})

\tkzFctPar[smooth,samples=250,domain=0:22.5147,->,>=latex,line width=0.8pt,color=red]{(0.3+6*t/225.147)*cos(t)}{(0.3+6*t/225.147)*sin(t)} 
\tkzDefShiftPoint[O]({\a*180/3.14159}:1){P}
\tkzDrawPoints(O,P) 
\tkzDefShiftPoint[O]({\a*180/3.14159}:1.6){B}
\tkzLabelPoint[below](B){$B$}
\tkzDefShiftPoint[O]({\a*90/3.14159}:1){T}	
\tkzLabelPoint[above left](T){\small\color{red}$\tfrac{7\pi}{6}+6\pi~\si{\rad}$}	 
\tkzDrawLine[line width=1.2pt,add=.0 and .6,->,>=latex,color=red](O,P)
%		\tkzDrawLine[line width=1.2pt,add=.0 and .4,->,>=latex](O,I)
\end{tikzpicture} 
\caption{$(\protect\vv{OA}~;~\protect\vv{OB})=-\frac{5\pi}{6}+2k\pi$}
\end{figure}  


\begin{example}[déterminer la valeur principale d'un angle]\hfill \\
Retrouver les valeurs principales des angles $\tfrac{7\pi}{2}$ et $-\tfrac{7\pi}{4}$.
\end{example}
\begin{proof}[solution]   $\begin{WithArrows}
		3\pi<\tfrac{7\pi}{2}&\leqslant \tikz[baseline =(A.base),remember picture]\node[draw,ellipse, color=red, line width=2pt](A) {$4\pi$};\Arrow{$-4\pi$}\\
		-\pi <\tfrac{7\pi}{2}-4\pi &\leqslant 0\\
		-\pi <-\tfrac{\pi}{2} &\leqslant 0
		\end{WithArrows}$ \hfill 
	$\begin{WithArrows}
		\tikz[baseline =(A.base),remember picture]\node[draw,ellipse, color=red, line width=2pt](A) {$-2\pi$}; <-\tfrac{7\pi}{4}&\leqslant -\pi\Arrow{$+2\pi$}\\
		0 <-\tfrac{7\pi}{4}+2\pi &\leqslant \pi\\
		0 <\tfrac{\pi}{4} &\leqslant \pi
	\end{WithArrows}$

Les mesures principales sont $-\tfrac{\pi}{2}$ et $\tfrac{\pi}{4}$
\end{proof}
%----------------------------------------------------------------------------------------
%	 
%----------------------------------------------------------------------------------------
\newpage 
\loadgeometry{widelayout}  
\pagestyle{scrheadings} % Print headers again  
\KOMAoptions{
	headwidth=textwithmarginpar,
	headsepline=true,
	headsepline= 0.5pt,
	footwidth=textwithmarginpar,
}  
\onehalfspacing 
\subsection{Exercices angles et radians}



\begin{exercise}[concepts] Complétez
	\begin{spacing}{2}
		\begin{enumerate}[label={},leftmargin=*]
			\item La mesure en radian d'un angle $\theta$ est \dotfill de l'arc intercepté par l'angle sur un cercle de rayon \dotfill
			\item Pour convertir en radians il faut multiplier les mesures données en degrés par \dotfill 
			\item Si l'angle au centre est $\theta$ et le rayon du cercle $r$, alors la longueur de l'arc intercepté est \dotfill 
			\item La mesure principale d'un angle orienté est toujours comprise entre \dotfill 
			\item   $(\vv{OI}~;~\vv{OP})=\theta+2k\pi$. Si $\theta\in\interval{\tfrac{\pi}{2}}{\pi}$ alors $P\in$ Quadrant \dotfill 
		\end{enumerate}
	\end{spacing} \vspace{-1em}
\end{exercise}  
\begin{exercise} Convertir en degrés les mesures données en~\si{\rad} 
	\vspace{-1em}
	\begin{spacing}{1.8}
		\begin{multicols}{5} 
			\begin{enumerate}[label={$\theta_{\arabic*}=$},leftmargin=*]
				\item  $\tfrac{5\pi}{3}~\si{\rad}$
				\item $\tfrac{3\pi}{4}~\si{\rad}$
				\item $\tfrac{5\pi}{6}~\si{\rad}$
				\item $3~\si{\rad}$
				\item $2~\si{\rad}$
			\end{enumerate}
		\end{multicols}\vspace{-2.5em}
	\end{spacing}
\end{exercise}
\begin{exercise} Convertir en~\si{\rad}  les mesures données en degrés
	\vspace{-1em}
\begin{spacing}{1.8}
	\begin{multicols}{5} 
		\begin{enumerate}[label={$\theta_{\arabic*}=$},leftmargin=*]
			\item  \ang{15}
			\item  \ang{36}
			\item  \ang{54}
			\item  \ang{75}
			\item  \ang{3600}
		\end{enumerate}
	\end{multicols}\vspace{-2.5em}
\end{spacing}
\end{exercise}
\begin{exercise}\hfill \vspace{-0.5em}
\begin{multicols}{2}
\begin{enumerate}[leftmargin=*,label=\arabic*)]
	\item Dans chaque cas déterminer la longueur $s$ de l'arc\\
	\begin{tikzpicture}[scale=1.8];
		\tkzSetUpPoint[size=3, fill=black,] %shape=cross out
		\tkzfctset{tan style/.style={->,>=latex, color=red, line width=1.2pt}}   
		\tkzSetUpLine[line width= 1.2pt, add=.2 and .2] 
		%		
		\tkzDefPoint(0,0){O}  
		\tkzDefPoint(1,0){I}
		\tkzDefShiftPoint[O](-30:1){A}
		\tkzDefShiftPoint[O](120:1){B}
		\tkzFillAngle[size=1.0,left color=gray!70,
		right color=white](A,O,B)
		
		\tkzDrawArc[angles,color=blue,line width=1.2pt](O,I)(120,-30)
		\tkzDrawArc[angles,color=red, line width=1.2pt, -,>=latex](O,I)(-30,120)
		\tkzDrawSegments[line width=1pt,-,>=latex](O,A O,B)
		\tkzMarkAngle[size=0.4](A,O,B)
		\tkzLabelAngle[pos=0.5](A,O,B){$\tfrac{5\pi}{6}$}
		\tkzLabelSegment[  below, font=\footnotesize](O,A'){$9$}
		%		\tkzLabelPoints[right](A,B)
		\tkzLabelPoints[left](O)
		\tkzText[text width=2.5cm](1.8,0.5){$s$}
	\begin{scope}[xshift=2.5cm] 
		\tkzSetUpPoint[size=3, fill=black,] %shape=cross out
		\tkzfctset{tan style/.style={->,>=latex, color=red, line width=1.2pt}}   
		\tkzSetUpLine[line width= 1.2pt, add=.2 and .2] 
		%		
		\tkzDefPoint(0,0){O}  
		\tkzDefPoint(1,0){I}
		\tkzDefShiftPoint[O](0:1){A}
		\tkzDefShiftPoint[O](140:1){B}
		\tkzFillAngle[size=1.0,left color=gray!70,
		right color=white](B,O,A)
		
		\tkzDrawArc[angles,color=red,line width=1.2pt](O,I)(140,0)
		\tkzDrawArc[angles,color=blue, line width=1.2pt, -,>=latex](O,I)(0,140)
		\tkzDrawSegments[line width=1pt,-,>=latex](O,A O,B)
		\tkzMarkAngle[size=0.4](A,O,B)
		\tkzLabelAngle[pos=0.5](A,O,B){$\ang{140}$}
		\tkzLabelSegment[  below, font=\footnotesize](O,A){$5$}
		%		\tkzLabelPoints[right](A,B)
		\tkzLabelPoints[left](O)
		\tkzText[text width=2.5cm](-0.3,-0.75){$s$}
		\end{scope}
	\end{tikzpicture}  
\item Déterminer la mesure de $\theta$ en~\si{\rad} puis en degrés :\\
	\begin{tikzpicture}[scale=1.8];
	\tkzSetUpPoint[size=3, fill=black,] %shape=cross out
	\tkzfctset{tan style/.style={->,>=latex, color=red, line width=1.2pt}}   
	\tkzSetUpLine[line width= 1.2pt, add=.2 and .2] 
	%		
	\tkzDefPoint(0,0){O}  
	\tkzDefPoint(1,0){I}
	\tkzDefShiftPoint[O](-30:1){A}
	\tkzDefShiftPoint[O](90:1){B} 
	
	\tkzDrawArc[angles,color=blue,line width=1.2pt](O,I)(90,-30)
	\tkzDrawArc[angles,color=red, line width=1.2pt, -,>=latex](O,I)(-30,90)
	\tkzDrawSegments[line width=1pt,-,>=latex](O,A O,B)
	\tkzMarkAngle[size=0.4](A,O,B)
	\tkzLabelAngle[pos=0.5](A,O,B){$\theta$}
	\tkzLabelSegment[  below, font=\footnotesize](O,A'){$4$}
	%		\tkzLabelPoints[right](A,B)
	\tkzLabelPoints[left](O)
	\tkzText[text width=2.5cm](1.8,0.5){$8$}
	\begin{scope}[xshift=2.5cm] 
		\tkzSetUpPoint[size=3, fill=black,] %shape=cross out
		\tkzfctset{tan style/.style={->,>=latex, color=red, line width=1.2pt}}   
		\tkzSetUpLine[line width= 1.2pt, add=.2 and .2] 
		%		
		\tkzDefPoint(0,0){O}  
		\tkzDefPoint(1,0){I}
		\tkzDefShiftPoint[O](0:1){A}
		\tkzDefShiftPoint[O](140:1){B} 
		
		\tkzDrawArc[angles,color=blue,line width=1.2pt](O,I)(140,0)
		\tkzDrawArc[angles,color=red, line width=1.2pt, -,>=latex](O,I)(0,140)
		\tkzDrawSegments[line width=1pt,-,>=latex](O,A O,B)
		\tkzMarkAngle[size=0.4](A,O,B)
		\tkzLabelAngle[pos=0.5](A,O,B){$\theta$}
		\tkzLabelSegment[  below, font=\footnotesize](O,A){$12$}
		%		\tkzLabelPoints[right](A,B)
		\tkzLabelPoints[below](O)
		\tkzText[text width=2.5cm](1.6,0.75){$16$}
	\end{scope}
\end{tikzpicture}  
\end{enumerate}
\end{multicols}
\end{exercise}
\begin{exercise} $\mathscr{C}$ est le cercle unité dans le repère $(O;~I,J)$ Placer les points $P_i$ sur $\mathscr{C}$ sachant que :\\
\noindent
\parbox{0.25\linewidth}{
\begin{tikzpicture}[scale=2]
	\tkzSetUpPoint[size=3, fill=black]%,shape=cross out
	\tkzfctset{tan style/.style={->,>=latex, color=red, line width=1.2pt}}   
	\tkzSetUpLine[line width= 1.2pt, add=.2 and .2]  
	\tkzInit[xmin=-1.0,xmax=1.0,ymin=-1.0,ymax=1.0,xstep=1.0, ystep=1.0];   
	\tkzDefPoint(0,0){O}  
	\tkzDefPoint(1,0){I}   
	\tkzDefPoint(0,1){J}  
	\tkzDrawX[ font=\small, line width=1.2pt, label={}  ,below left, left space=0.1,right space=0.1] 
	\tkzDrawY[ font=\small, line width=1.2pt , label={} ,below left, down space=0.1,up space=0.1]  
	\tkzDrawCircle[line width=2pt,color=black](O,I)
	\foreach \x/\xtext in {  
		0/I,
		90/J,
		30/A 
	} 
	{
		\tkzDefShiftPoint[O](\x:1){\xtext};
		\tkzDrawPoints(\xtext);  
		\tkzDrawSegment[line width=1pt](O,\xtext); 
		\tkzText[fill=white](\x:1.20){\scriptsize$\xtext$};
	} 
\tkzDefShiftPoint[O](30:1){A};
\tkzMarkAngle[size=0.3](I,O,A) 
\tkzLabelAngle[pos=0.5](I,O,A){\scriptsize\ang{30}}
\foreach \x in {
	60,
	120,
	150,
	210,
	240,
	300,
	330 ,
	45,
	135,
	-45,
	-135
} 
{
	\tkzDefShiftPoint[O](\x:1){B};
	\tkzDrawPoints(B);  
	\tkzDrawSegment[line width=0.5pt,dashed](O,B);  
}  
\end{tikzpicture} 
}
\hfill
\parbox{0.35\linewidth}{	
$(\vv{OI}~;~\vv{OP_1})=-\tfrac{5\pi}{6}$\\   
$(\vv{OA}~;~\vv{OP_2})=\tfrac{\pi}{2}-6\pi$\\   
$(\vv{OI}~;~\vv{OP_3})=\pi - (\vv{OI}~;~\vv{OA})$\\  
$(\vv{OI}~;~\vv{OP_4})=\pi + (\vv{OI}~;~\vv{OA})$   
}
\parbox{0.35\linewidth}{	 
	$(\vv{OJ}~;~\vv{OP_5})=\tfrac{7\pi}{4}+10\pi$\\   
	$(\vv{OA}~;~\vv{OP_6})=2023\pi$\\   
	$(\vv{OA}~;~\vv{OP_7})=-\tfrac{3\pi}{2}+5\pi$\\   
	$(\vv{OI}~;~\vv{OP_8})=-\tfrac{5\pi}{4}+7\pi$
}	
\end{exercise}
\begin{exercise} Déterminer la mesure principale des angles orientés donnés en~\si{\rad}.
	\vspace{-1em}
	\begin{spacing}{1.8}
		\begin{multicols}{5} 
			\begin{enumerate}[label={$\theta_{\arabic*}=$},leftmargin=*]
				\item $173\pi$
				\item $-250\pi$
				\item $\tfrac{7\pi}{3}$
				\item $-\tfrac{17\pi}{6}$
				\item $\tfrac{53\pi}{2}$
			\end{enumerate}
		\end{multicols}\vspace{-2.5em}
	\end{spacing} 
\end{exercise}
\begin{exercise} $P$ est un point du plan repéré. Dans chaque cas déterminez à quel quadrant il appartient.
	\vspace{-1em}
	\begin{spacing}{1.8}
		\begin{multicols}{4} 
			\begin{enumerate}[label={\arabic*)},leftmargin=*] 
				\item  $(\vv{OI}~;~\vv{OP})=\tfrac{21\pi}{2}$
				\item  $(\vv{OI}~;~\vv{OP})=\tfrac{14\pi}{3}$
				\item  $(\vv{OI}~;~\vv{OP})=\tfrac{-13\pi}{4}$
				\item  $(\vv{OI}~;~\vv{OP})=\tfrac{5\pi}{3}-2\pi$
%				\item  $(\vv{OI}~;~\vv{OP})=\tfrac{7\pi}{6}-3\pi$
			\end{enumerate}
		\end{multicols}\vspace{-2.5em}
	\end{spacing} 
\end{exercise}

%----------------------------------------------------------------------------------------
%	 
%----------------------------------------------------------------------------------------
\newpage 
\loadgeometry{marginlayout}  
\pagestyle{scrheadings} % Print headers again  
\KOMAoptions{
	headwidth=textwithmarginpar,
	headsepline=true,
	headsepline= 0.5pt,
	footwidth=textwithmarginpar,
}  
\onehalfspacing 
\section{Fonctions trigonométriques} 
\begin{marginfigure}
	\begin{tikzpicture}[scale=2]
		\tkzSetUpPoint[size=3, fill=black]%,shape=cross out
		\tkzfctset{tan style/.style={->,>=latex, color=red, line width=1.2pt}}   
		\tkzSetUpLine[line width= 1.2pt, add=.2 and .2]  
		\tkzInit[xmin=-1.0,xmax=1.0,ymin=-1.0,ymax=1.0,xstep=1.0, ystep=1.0];  
		
		\tkzDrawX[ font=\small, line width=2pt, label=$x$  ,left space=0.4,right space=0.4] 
		\tkzDrawY[ font=\small, line width=2pt , label={$y$} , down space=0.3,up space=0.3]  
		%	\tkzLabelXY[orig = false,  font=\scriptsize ] 
		%	\tkzRep[color  = Maroon,xlabel=,ylabel=, font=\scriptsize, line width = 2pt,xnorm=1, ynorm=1]   
		\def\a{55} 
		\tkzDefPoint(0,0){O}  
		\tkzDefPoint(1,0){I}  
		\tkzDefPoint(0,1){J}
		\tkzLabelPoint[below left](O){$0$}
		\tkzDrawCircle[color=blue, line width=1.2pt](O,I)
		\tkzLabelPoint[below right](I){$1$}
		\tkzLabelPoint[above right](J){$1$}
		\tkzDrawArc[color=red,->,>=latex,angles,line width=2pt](O,I)(0,\a)
		\tkzDefShiftPoint[O](\a:1){P}
		\tkzDrawPoints(P)
		\tkzLabelPoint[above right](P){$P(x~;~y)$}	
		\tkzDefShiftPoint[O]({\a/2}:1){T}	
		\tkzLabelPoint[above right](T){\color{red}$t$}	
		%		\tkzDefShiftPoint[O](70:3){A}
		%		\tkzDefShiftPoint[O](-90:3){B} 
		%		\tkzDrawPoints(A,B)
		%		\tkzLabelPoint[above right, font=\scriptsize](A){$A(2;y)$}
		%		\tkzLabelPoint[below right, font=\scriptsize](B){$B(0;9)$}  
		\tkzPointShowCoord[xlabel=$\cos{t}$,ylabel=$\sin{t}$](P)
		\tkzDrawSegment[dashed,](O,P)
	\end{tikzpicture} 
\caption{Définition de $\cos$ et $\sin$}
\end{marginfigure}
\begin{definition} Pour tout $t\in\mathbb{R}$,  $P(x~;~y)$ désigne l'unique point correspondant à $t$ sur le cercle unité :
	\begin{enumerate}[label=\textbullet,leftmargin=*]
		\item Le \textbf{cosinus} de $t$, est l'abscisse de $P$ \hfill  $\cos{t}=x$
		\item Le \textbf{sinus} de $t$, est l'ordonnée de $P$ \hfill $\sin{t}=x$
	\end{enumerate}
	Les coordonnées de $P$ sont : $P(\cos(t)~;~\sin(t))$.
\end{definition}
\begin{property} Pour tout $t\in\mathbb{R}$ : $-1\leqslant \cos t~;~\sin t \leqslant 1$ et 
	\setlength{\abovedisplayskip}{3pt}
	\setlength{\belowdisplayskip}{3pt} 
	$$
	\text{pour tout $t\in\mathbb{R}$}\qquad
	\cos^2(t) + \sin^2(t) = 1
	$$
\end{property} 
\begin{property}  Les fonctions $\cos$ et $\sin$ sont périodiques de \textbf{période} $2\pi$ : 
	\setlength{\abovedisplayskip}{3pt}
	\setlength{\belowdisplayskip}{3pt}  
	\begin{align*}
		 \text{pour tout $t\in\mathbb{R}$ et $n\in\mathbb{Z}$}\qquad
		 \cos(t+2n\pi)&=\cos(t)\\ 
		 \sin(t+2n\pi)&=\sin(t) 
	\end{align*} 
\end{property}
\begin{property}[parité] 
\begin{marginfigure} 
	\begin{tikzpicture}[scale=2]
		\tkzSetUpPoint[size=3, fill=black]%,shape=cross out
		\tkzfctset{tan style/.style={->,>=latex, color=red, line width=1.2pt}}   
		\tkzSetUpLine[line width= 1.2pt, add=.2 and .2]  
		\tkzInit[xmin=-1.0,xmax=1.0,ymin=-1.0,ymax=1.0,xstep=1.0, ystep=1.0];  
		
		
		%	\tkzLabelXY[orig = false,  font=\scriptsize ] 
		%	\tkzRep[color  = Maroon,xlabel=,ylabel=, font=\scriptsize, line width = 2pt,xnorm=1, ynorm=1]   
		\def\a{40} 
		\def\b{-40} 
		\def\c{180-40} 
		\def\d{180+40} 
		\tkzDefPoint(0,0){O}  
		\tkzDefPoint(1,0){I}  
		\tkzDefPoint(-1,0){I'}  
		\tkzDefPoint(0,1){J}
		\tkzLabelPoint[below left](O){$0$}
		\tkzDrawCircle[color=blue, line width=1.2pt](O,I)
		\tkzLabelPoint[below right](I){$1$}
		\tkzLabelPoint[above right](J){$1$}
		\tkzDrawArc[color=red,->,>=latex,angles,line width=2pt](O,I)(0,\a)
		\tkzDrawArc[color=red,<-,>=latex,angles,line width=2pt](O,I)(\b,0)
		\tkzDrawArc[color=red,->,>=latex,angles,line width=2pt](O,I)(180,\d)
		\tkzDrawArc[color=red,<-,>=latex,angles,line width=2pt](O,I)(\c,180)
		\tkzDefShiftPoint[O](\a:1){P}
		\tkzDrawPoints(I,I')
		\tkzDrawPoints(P)
		\tkzDefShiftPoint[O](\b:1){P'}
		\tkzDrawPoints(P')
		\tkzDefShiftPoint[O](\d:1){Q}
		\tkzDrawPoints(Q)
		\tkzDefShiftPoint[O](\c:1){Q'}
		\tkzDrawPoints(Q')
		\tkzLabelPoint[above right,font=\scriptsize](P){$P(t)$}	
		\tkzLabelPoint[below right,font=\scriptsize](P'){$P(-t)$}
		\tkzLabelPoint[below ,rotate=\b ,font=\scriptsize](Q){$P(\pi+t)$}
		\tkzLabelPoint[above ,rotate=\a ,font=\scriptsize](Q'){$P(\pi-t)$}
		\tkzDefShiftPoint[O]({\a/2}:1){T}	
		\tkzLabelPoint[ right,font=\scriptsize](T){\color{red}$t$}	
		\tkzDefShiftPoint[O]({\b/2}:1){T'}	
		\tkzLabelPoint[ right,font=\scriptsize](T'){\color{red}$-t$}
		\tkzDefShiftPoint[O]({-\c/2}:1){T0}	
		\tkzLabelPoint[ left,font=\scriptsize](T0){\color{red}$-t$}	
		\tkzDefShiftPoint[O]({\d/2}:1){T1}	
		\tkzLabelPoint[ left,font=\scriptsize](T1){\color{red}$t$}	
		%		\tkzDefShiftPoint[O](70:3){A}
		%		\tkzDefShiftPoint[O](-90:3){B} 
		%		\tkzDrawPoints(A,B)
		%		\tkzLabelPoint[above right, font=\scriptsize](A){$A(2;y)$}
		%		\tkzLabelPoint[below right, font=\scriptsize](B){$B(0;9)$}  
		\tkzPointShowCoord[xlabel=$\cos{t}$,  ylabel=$\sin{t}$](P)
%		\tkzPointShowCoord[ ylabel=$-\sin{t}$](P')
%		\tkzDrawSegments[dashed](O,P O,P')
		\tkzDrawX[ font=\small, line width=2pt, label=$x$  ,left space=0.4,right space=0.4] 
		\tkzDrawY[ font=\small, line width=2pt , label={$y$} , down space=0.3,up space=0.3]   
	\end{tikzpicture} 
	\caption{Parité}
\end{marginfigure} 
Pour tout $t\in\mathbb{R}$
\setlength{\abovedisplayskip}{3pt}
\setlength{\belowdisplayskip}{3pt}  
\begin{align*} 
	\cos(-t)&=\cos(t)&&\text{la fonction $\cos$ est paire}\\ 
	\sin(-t)&=-\sin(t)&&\text{la fonction $\sin$ est impaire}  
\end{align*}  
\end{property} 
\begin{property} Pour tout $t\in\mathbb{R}$ :
	\setlength{\abovedisplayskip}{3pt}
	\setlength{\belowdisplayskip}{3pt}  
	\begin{align*}
	\cos(\pi+t)&=-\cos(t) &  	\sin(\pi+t)&=-\sin(t) \\
	\cos(\pi-t)&=-\cos(t) &  	\sin(\pi-t)&=\sin(t)  
	\end{align*}
\end{property} 
\begin{remark}
$k\in\interval{-1}{1}$ a une infinité d'antécédents par    $\cos$ ou $\sin$. 
\begin{enumerate}[leftmargin=*,label=\textbullet]
	\item $\arccos(k)$ est l'antécédent de $k$ par $\cos$ sur $\interval{0}{\pi}$.
	\item $\arcsin(k)$ est l'antécédent de $k$ par $\sin$ sur $\interval{-\tfrac{\pi}{2}}{\tfrac{\pi}{2}}$.
\end{enumerate}
\end{remark}
\begin{figure*}[!h]
		\begin{tikzpicture}[yscale=0.9,xscale=1.5]
		\tkzSetUpPoint[size=3, fill=black]
		\tkzfctset{tan style/.style={->,>=latex, color=red, line width=1.2pt}}   
		\tkzSetUpLine[line width= 1.2pt, add=.5 and .5]  
		\tkzInit[xmin=-5,xmax=7,ymin=-1.0,ymax=1,xstep=1, ystep=0.5]; 
		%	\begin{scope}[dashed]  
%					\tkzGrid[   color=black]  
			%	\end{scope}
		\tkzDrawX[above=5pt, label={$x$}, font=\small, line width=1.2pt, trig=2] 
		\tkzDrawY[right=5pt, label={$y$}, font=\small, up space =0.25, line width=1.2pt]
		%	\tkzText[font=\small](60,-20){temps en s}
		%	\tkzText[rotate=90,font=\small](-15,100){distance à la maison en m}
		\tkzLabelX[  font=\scriptsize , trig=2]
		\tkzLabelY[orig=false,  font=\scriptsize]    
			\tkzFct[domain =-2*pi:6*pi/2,  line width=0.5pt,dashed]{ 
			(cos(x))	
		};   
		\tkzFct[domain =0:pi,  line width=2pt,color=red]{ 
			(cos(x))	
		}; 
%\tkzVTicks{0.5,-0.5,1,-1,sqrt(3)/2}  
\foreach\v in {sqrt(2)/2}
{\tkzHLine[color=blue,line width=0.5pt,dashed]{\v}}
\tkzText[font=\large,draw, fill=white](pi+0.5,1){$\cos$}
		%		\tkzSetUpPoint[size=3, fill=black]
%		\tkzDrawTangentLine[kr=0.2,kl=0.2,draw](0)
%		\tkzDrawTangentLine[kr=0.5,kl=0.5,draw](1)
%		\tkzDrawTangentLine[kr=0.5,kl=0.5,draw](4)
		%		\tkzDrawTangentLine[kr=0.75,kl=0.75,draw](1)
%		\tkzDefPointByFct[ref=A]({0}); 
		%		\tkzDefPointByFct[ref=B]({1});  
		%		\tkzLabelPoint[below](A){\small$A$}
		%		\tkzLabelPoint[above](B){\small$B$}
%		\tkzVLine[dashed, line width=2pt]{1}	
%		\tkzText[text width=5cm](3,1){$f$ est croissante sur $\interval[open right]{1}{+\infty}$. la dérivée $f'(x)\geqslant 0$.} 
%		\tkzText[text width=5.5cm](-2,-1.5){$f$ est décroissante sur $\interval[open left]{-\infty}{1}$. la dérivée $f'(x)\leqslant 0$.}
	\end{tikzpicture}
		\begin{tikzpicture}[yscale=1.8,xscale=1.5]
\tkzSetUpPoint[size=3, fill=black]
\tkzfctset{tan style/.style={->,>=latex, color=red, line width=1.2pt}}   
\tkzSetUpLine[line width= 1.2pt, add=.5 and .5]  
\tkzInit[xmin=-5,xmax=7,ymin=-1.0,ymax=1,xstep=1.0, ystep=1.0]; 
%	\begin{scope}[dashed]  
	%		\tkzGrid[   color=gray](-6,-6)(7,5); 
	%	\end{scope}
\tkzDrawX[above=5pt, label={$x$}, font=\small, line width=1.2pt, trig=2] 
\tkzDrawY[right=5pt, label={$y$}, font=\small,up space=0.36, line width=1.2pt]
%	\tkzText[font=\small](60,-20){temps en s}
%	\tkzText[rotate=90,font=\small](-15,100){distance à la maison en m}
\tkzLabelX[  font=\scriptsize , trig=2]
\tkzLabelY[orig=false,  font=\scriptsize]    
\tkzFct[domain =-2*pi:6*pi/2,  line width=1pt,dashed]{ 
	(sin(x))	
};   
\tkzFct[domain =-pi/2:pi/2,  line width=2pt,color=red]{ 
	(sin(x))	
};  
%\tkzVTicks{0.5,-0.5,1,-1,sqrt(3)/2}  
\foreach\v in {sqrt(2)/2}
{\tkzHLine[color=blue,line width=0.5pt,dashed]{\v}}
\tkzText[font=\large,draw, fill=white](pi+0.5,1){$\sin$}
%		\tkzSetUpPoint[size=3, fill=black]
%		\tkzDrawTangentLine[kr=0.2,kl=0.2,draw](0)
%		\tkzDrawTangentLine[kr=0.5,kl=0.5,draw](1)
%		\tkzDrawTangentLine[kr=0.5,kl=0.5,draw](4)
%		\tkzDrawTangentLine[kr=0.75,kl=0.75,draw](1)
%		\tkzDefPointByFct[ref=A]({0}); 
%		\tkzDefPointByFct[ref=B]({1});  
%		\tkzLabelPoint[below](A){\small$A$}
%		\tkzLabelPoint[above](B){\small$B$}
%		\tkzVLine[dashed, line width=2pt]{1}	
%		\tkzText[text width=5cm](3,1){$f$ est croissante sur $\interval[open right]{1}{+\infty}$. la dérivée $f'(x)\geqslant 0$.} 
%		\tkzText[text width=5.5cm](-2,-1.5){$f$ est décroissante sur $\interval[open left]{-\infty}{1}$. la dérivée $f'(x)\leqslant 0$.}
\end{tikzpicture}
\caption{Représentations graphiques des fonctions $\cos$ et $\sin$}
\end{figure*} 
\begin{property}[Équation $\cos{t}= \cos{a}$ et $\sin{t}= \sin{a}$]
	Soit $t$ et $a\in\mathbb{R}$ :
\begin{enumerate}[label=\textbullet,leftmargin=*]
	\item $\cos(t)=\cos(a) \iff t= a+2k\pi  \quad\text{ou}\quad t= -a +2k\pi , \quad k\in\mathbb{Z}$\\
	\noindent
	\parbox{1\linewidth}{
	\centering
		\begin{tikzpicture}[scale=2]
		\tkzSetUpPoint[size=3, fill=black]%,shape=cross out
		\tkzfctset{tan style/.style={->,>=latex, color=red, line width=1.2pt}}   
		\tkzSetUpLine[line width= 1.2pt, add=.2 and .2]  
		\tkzInit[xmin=-1.0,xmax=1.0,ymin=-1.0,ymax=1.0,xstep=1.0, ystep=1.0];  
		
		
		%	\tkzLabelXY[orig = false,  font=\scriptsize ] 
		%	\tkzRep[color  = Maroon,xlabel=,ylabel=, font=\scriptsize, line width = 2pt,xnorm=1, ynorm=1]   
		\def\a{50} 
		\def\b{-50} 
		\def\c{180-50} 
		\def\d{180+50} 
		\tkzDefPoint(0,0){O}  
		\tkzDefPoint(1,0){I}  
		\tkzDefPoint(-1,0){I'}  
		\tkzDefPoint(0,1){J}
		\tkzLabelPoint[below left](O){$0$}
		\tkzDrawCircle[color=blue, line width=1.2pt](O,I)
		\tkzLabelPoint[below right](I){$1$}
		\tkzLabelPoint[above right](J){$1$}
		\tkzDrawArc[color=red,->,>=latex,angles,line width=2pt](O,I)(0,\a)
		\tkzDrawArc[color=red,<-,>=latex,angles,line width=2pt](O,I)(\b,0)
%		\tkzDrawArc[color=red,->,>=latex,angles,line width=2pt](O,I)(180,\d)
%		\tkzDrawArc[color=red,<-,>=latex,angles,line width=2pt](O,I)(\c,180)
		\tkzDefShiftPoint[O](\a:1){P}
		\tkzDrawPoints(I,I')
		\tkzDrawPoints(P)
		\tkzDefShiftPoint[O](\b:1){P'}
		\tkzDrawPoints(P')
%		\tkzDefShiftPoint[O](\d:1){Q}
%		\tkzDrawPoints(Q)
%		\tkzDefShiftPoint[O](\c:1){Q'}
%		\tkzDrawPoints(Q')
		\tkzLabelPoint[above right,font=\scriptsize](P){$P(t)$}	
		\tkzLabelPoint[below right,font=\scriptsize](P'){$P(-t)$}
%		\tkzLabelPoint[below ,rotate=\b ,font=\scriptsize](Q){$P(\pi+t)$}
%		\tkzLabelPoint[above ,rotate=\a ,font=\scriptsize](Q'){$P(\pi-t)$}
		\tkzDefShiftPoint[O]({\a/2}:1){T}	
		\tkzLabelPoint[ right,font=\scriptsize](T){\color{red}$t$}	
		\tkzDefShiftPoint[O]({\b/2}:1){T'}	
		\tkzLabelPoint[ right,font=\scriptsize](T'){\color{red}$-t$}
%		\tkzDefShiftPoint[O]({-\c/2}:1){T0}	
%		\tkzLabelPoint[ left,font=\scriptsize](T0){\color{red}$-t$}	
%		\tkzDefShiftPoint[O]({\d/2}:1){T1}	
%		\tkzLabelPoint[ left,font=\scriptsize](T1){\color{red}$t$}	
		%		\tkzDefShiftPoint[O](70:3){A}
		%		\tkzDefShiftPoint[O](-90:3){B} 
		%		\tkzDrawPoints(A,B)
		%		\tkzLabelPoint[above right, font=\scriptsize](A){$A(2;y)$}
		%		\tkzLabelPoint[below right, font=\scriptsize](B){$B(0;9)$}   
		\tkzVLine[color=blue,line width=0.5pt,dashed]{0.6427}
		\tkzPointShowCoord[noydraw,xstyle={below=10pt} ,  xlabel=$\cos{t}$](P)
		%		\tkzPointShowCoord[ ylabel=$-\sin{t}$](P')
		%		\tkzDrawSegments[dashed](O,P O,P')
\tkzDrawX[ font=\small, line width=2pt, label=$x$  ,left space=0.4,right space=0.4] 
\tkzDrawY[ font=\small, line width=2pt , label={$y$} , down space=0.3,up space=0.3]  
	\end{tikzpicture} 
\hfill 
	\begin{tikzpicture}[scale=2]
	\tkzSetUpPoint[size=3, fill=black]%,shape=cross out
	\tkzfctset{tan style/.style={->,>=latex, color=red, line width=1.2pt}}   
	\tkzSetUpLine[line width= 1.2pt, add=.2 and .2]  
	\tkzInit[xmin=-1.0,xmax=1.0,ymin=-1.0,ymax=1.0,xstep=1.0, ystep=1.0];  
	
	
	%	\tkzLabelXY[orig = false,  font=\scriptsize ] 
	%	\tkzRep[color  = Maroon,xlabel=,ylabel=, font=\scriptsize, line width = 2pt,xnorm=1, ynorm=1]   
	\def\a{40} 
	\def\b{-40} 
	\def\c{180-40} 
	\def\d{180+40} 
	\tkzDefPoint(0,0){O}  
	\tkzDefPoint(1,0){I}  
	\tkzDefPoint(-1,0){I'}  
	\tkzDefPoint(0,1){J}
	\tkzLabelPoint[below left](O){$0$}
	\tkzDrawCircle[color=blue, line width=1.2pt](O,I)
	\tkzLabelPoint[below right](I){$1$}
	\tkzLabelPoint[above right](J){$1$}
	\tkzDrawArc[color=red,->,>=latex,angles,line width=2pt](O,I)(0,\a)
%	\tkzDrawArc[color=red,<-,>=latex,angles,line width=2pt](O,I)(\b,0)
	%		\tkzDrawArc[color=red,->,>=latex,angles,line width=2pt](O,I)(180,\d)
			\tkzDrawArc[color=red,<-,>=latex,angles,line width=2pt](O,I)(\c,180)
	\tkzDefShiftPoint[O](\a:1){P}
	\tkzDrawPoints(I,I')
	\tkzDrawPoints(P)
%	\tkzDefShiftPoint[O](\b:1){P'}
%	\tkzDrawPoints(P')
	%		\tkzDefShiftPoint[O](\d:1){Q}
	%		\tkzDrawPoints(Q)
	%		\tkzDefShiftPoint[O](\c:1){Q'}
			\tkzDrawPoints(Q')
	\tkzLabelPoint[above right,font=\scriptsize](P){$P(t)$}	
%	\tkzLabelPoint[below right,font=\scriptsize](P'){$P(-t)$}
	%		\tkzLabelPoint[below ,rotate=\b ,font=\scriptsize](Q){$P(\pi+t)$}
	\tkzLabelPoint[above ,rotate=\a ,font=\scriptsize](Q'){$P(\pi-t)$}
	\tkzDefShiftPoint[O]({\a/2}:1){T}	
	\tkzLabelPoint[ right,font=\scriptsize](T){\color{red}$t$}	
%	\tkzDefShiftPoint[O]({\b/2}:1){T'}	
%	\tkzLabelPoint[ right,font=\scriptsize](T'){\color{red}$-t$}
			\tkzDefShiftPoint[O]({-\c/2}:1){T0}	
			\tkzLabelPoint[ left,font=\scriptsize](T0){\color{red}$-t$}	
	%		\tkzDefShiftPoint[O]({\d/2}:1){T1}	
	%		\tkzLabelPoint[ left,font=\scriptsize](T1){\color{red}$t$}	
	%		\tkzDefShiftPoint[O](70:3){A}
	%		\tkzDefShiftPoint[O](-90:3){B} 
	%		\tkzDrawPoints(A,B)
	%		\tkzLabelPoint[above right, font=\scriptsize](A){$A(2;y)$}
	%		\tkzLabelPoint[below right, font=\scriptsize](B){$B(0;9)$}  
	\tkzHLine[color=blue,line width=0.5pt,dashed]{0.6427}
	\tkzPointShowCoord[noxdraw, ylabel=$\sin{t}$](P)
	%		\tkzPointShowCoord[ ylabel=$-\sin{t}$](P')
	%		\tkzDrawSegments[dashed](O,P O,P')
	\tkzDrawX[ font=\small, line width=2pt, label=$x$  ,left space=0.4,right space=0.4] 
	\tkzDrawY[ font=\small, line width=2pt , label={$y$} , down space=0.3,up space=0.3]  
\end{tikzpicture}
}
	\item $\sin(t)=\sin(a) \iff t= a+2k\pi  \quad\text{ou}\quad t= \pi-a +2k\pi , \quad k\in\mathbb{Z}$
\end{enumerate}  
\end{property} 
\begin{example} Résoudre dans $\mathbb{R}$ l'équation $\sin(x)=-\tfrac{\sqrt{3}}{2}$
\end{example}
\begin{proof}[solution] $\begin{WithArrows}
		\sin(x)& =-\tfrac{\sqrt{3}}{2} \Arrow{on note que $\sin\big(-\tfrac{\pi}{3}\big)=-\tfrac{\sqrt{3}}{2}$}\\
		\iff \sin(x)& = \sin\big(-\tfrac{\pi}{3}\big) \\
		\iff x= -\tfrac{\pi}{3} +2k\pi \qquad \text{ou}\quad x& =\pi-\tfrac{\pi}{3} +2k'\pi
	\end{WithArrows}$\hfill \null
\end{proof}
\begin{figure*}[!hb]
	\noindent
\parbox{0.70\linewidth}{
\begin{tikzpicture}[scale=4.5]
	\tkzSetUpPoint[size=4, fill=black]%,shape=cross out
	\tkzfctset{tan style/.style={->,>=latex, color=red, line width=1.2pt}}   
	\tkzSetUpLine[line width= 1.2pt, add=.2 and .2]  
	\tkzInit[xmin=-1.0,xmax=1.0,ymin=-1.0,ymax=1.0,xstep=1.0, ystep=1.0];   
	\tkzDefPoint(0,0){O}  
	\tkzDefPoint(1,0){I}   
	\tkzDefPoint(0,1){J} 
	\begin{scope}[dashed]  
 		\clip (0,0) circle (1cm);
		\tkzHLines[color=blue,line width=0.5pt,dashed]{-sqrt(3)/2,-sqrt(2)/2,-0.5,0.5,sqrt(2)/2,sqrt(3)/2};
		\tkzVLines[color=blue,line width=0.5pt,dashed]{-sqrt(3)/2,-sqrt(2)/2,-0.5,0.5,sqrt(2)/2,sqrt(3)/2};
	\end{scope}
\tkzDrawX[ font=\small, line width=1.2pt, label={$\cos$}  ,below left, left space=0.1,right space=0.1] 
\tkzDrawY[ font=\small, line width=1.2pt , label={$\sin$} ,below left, down space=0.1,up space=0.1]  
	\tkzDrawCircle[line width=2pt,color=black](O,I)
	\foreach \x/\t/\xtext/\ytext in {
		0/0/1/0,
		30/\tfrac{\pi}{6}/\tfrac{\sqrt{3}}{2}/\tfrac{1}{2},
		45/\tfrac{\pi}{4}/\tfrac{\sqrt{2}}{2}/\tfrac{\sqrt{2}}{2},
		60/\tfrac{\pi}{3}/\tfrac{1}{2}/\tfrac{\sqrt{3}}{2},
		90/\tfrac{\pi}{2}/0/1, 
		135/\tfrac{3\pi}{4}/-\tfrac{\sqrt{2}}{2}/\tfrac{\sqrt{2}}{2},
		-135/-\tfrac{3\pi}{4}/-\tfrac{\sqrt{2}}{2}/-\tfrac{\sqrt{2}}{2},
		-45/-\tfrac{\pi}{4}/\tfrac{\sqrt{2}}{2}/-\tfrac{\sqrt{2}}{2},
		180/\pi/-1/0,
		-90/\tfrac{\pi}{2}/0/1
} 
	{
		\tkzDefShiftPoint[O](\x:1){P};
		\tkzDrawPoints(P);
		\tkzDefShiftPoint[O](\x:0.95){S};
		\tkzDefShiftPoint[O](\x:1.05){T};
		\tkzDrawSegment[line width=1.2pt](S,T);
		\tkzDrawSegment[line width=1pt](O,P);
		\tkzText[fill=white](\x:0.60){\scriptsize\ang{\x}};
		\tkzText[fill=white](\x:1.175){\scriptsize$\t$};
		\tkzText(\x:1.40){\small$\left(\xtext,\ytext\right)$};
	}
	\foreach \x/\t/\xtext/\ytext in { 
		120/\tfrac{2\pi}{3}/-\tfrac{1}{2}/\tfrac{\sqrt{3}}{2},
		150/\tfrac{5\pi}{6}/-\tfrac{\sqrt{3}}{2}/\tfrac{1}{2},
		-30/-\tfrac{\pi}{6}/\tfrac{\sqrt{3}}{2}/-\tfrac{1}{2},
		-60/-\tfrac{\pi}{3}/\tfrac{1}{2}/-\tfrac{\sqrt{3}}{2},
		-120/-\tfrac{2\pi}{3}/-\tfrac{1}{2}/-\tfrac{\sqrt{3}}{2},
		-150/-\tfrac{5\pi}{6}/-\tfrac{\sqrt{3}}{2}/-\tfrac{1}{2}
	} 
	{
		\tkzDefShiftPoint[O](\x:1){P};
		\tkzDefShiftPoint[O](\x:0.95){S};
		\tkzDefShiftPoint[O](\x:1.05){T};
		\tkzDrawSegment[line width=1.2pt](S,T);
		\tkzDefShiftPoint[O](\x:1){T};
		\tkzDrawPoints(P);
%		\tkzDrawSegment[line width=1pt](O,P);
		%		\tkzText[fill=white](\x:0.60){\scriptsize\ang{\x}};
		\tkzText[fill=white](\x:1.15){\scriptsize$\t$};
		\tkzText(\x:1.40){\small$\left(\xtext,\ytext\right)$};
	}
\end{tikzpicture} 
}
\parbox{0.25\linewidth}{
	\begin{tikzpicture}[scale=4]
		\draw[thick] (0.5,0) -- (0,0) -- (60:1) node[midway,above left]{1} --  (0.5,0)  node[midway,below right]{$\tfrac{\sqrt{3}}{2}$} -- (0,0) node[midway,below]{$\tfrac{1}{2}$}  ;
		\draw[dashed] (0.5,0) -- (1,0) -- (60:1)  ;
		\draw (0.5,0.1) -- ++(-0.1,0) -- ++(0,-0.1);
%		
		\draw (0.15,0) arc (0:60:0.15);
		\draw (30:0.15) node[right]{$\tfrac{\pi}{3}$};
%		
		\draw (0.5,{sqrt(3)/2-0.15}) arc (-90:-120:0.15);
		\draw ($(60:1) + (-105:0.3)$) node[]{$\tfrac{\pi}{6}$};
	\end{tikzpicture} 
	\begin{tikzpicture}[scale=3]
		\draw[thick] (0,0) -- (1,0) node[midway,below]{$\tfrac{\sqrt{2}}{2}$} -- (1,1) node[midway,right]{$\tfrac{\sqrt{2}}{2}$} -- (0,0)  node[midway,above left]{$1$};
		\draw (1,0.1) -- (0.9,0.1) -- (0.9,0);
		\draw (0.3,0) arc (0:45:0.3);
		\draw ($(0,0) + (22.5:0.3)$) node[right ]{$\tfrac{\pi}{4}$};
	\end{tikzpicture}\\
	\begin{tikzpicture}[scale=4]
		\draw[thick] (0,0) -- (30:1) node[midway,above left]{1} -- ++(0,-0.5) node[midway,right]{$\tfrac{1}{2}$} -- (0,0) node[midway,below]{$\tfrac{\sqrt{3}}{2}$}  ;
		\draw[dashed] (0,0) -- (-30:1) -- ++(0,0.5);
		\draw ({sqrt(3)/2},0.1) -- ++(-0.1,0) -- ++(0,-0.1);
		
		\draw (0.3,0) arc (0:30:0.3);
		\draw (15:0.3) node[right]{$\tfrac{\pi}{6}$};
		
		\draw ({sqrt(3)/2},0.3) arc (-90:-133:0.3);
		\draw ($(30:1) + (-110:0.3)$) node[]{$\tfrac{\pi}{3}$};
	\end{tikzpicture} 
} 
\caption{à mémoriser \emoji{heart} Coordonnées des points du cercle unité et valeurs particulières de $\cos$ et $\sin$}
 \end{figure*}

%----------------------------------------------------------------------------------------
%	 
%----------------------------------------------------------------------------------------
\newpage 
\loadgeometry{widelayout}  
\pagestyle{scrheadings} % Print headers again  
\KOMAoptions{
	headwidth=textwithmarginpar,
	headsepline=true,
	headsepline= 0.5pt,
	footwidth=textwithmarginpar,
}  
\onehalfspacing 
\subsection{Exercices fonctions trigonométriques} 
\begin{eBox}
	Dans ses exercices, le plan est muni d'un repère orthonormé $(O~;~I,J)$ et du cercle unité orienté $\mathscr{C}$.
\end{eBox}
\begin{exercise}[concepts] Complétez
\begin{spacing}{2}
\begin{enumerate}[label={\alph*)},leftmargin=*]
	\item  Pour $P(x~;~y)$ un point du cercle, et si $(\vv{OI}~;~\vv{OP})=t+2k\pi$ alors $\sin{t}=$\dotfill et $\cos{t}=$\dotfill
	\item Si $P(x~;~y)\in\mathscr{C}$ alors $x^2+y^2=$\dotfill. Donc $\cos^2 t +\sin^2 t=$\dotfill 
 \item Si $\sin(t)=0.6$, alors $\cos^2(t)=$\dotfill. Donc $\cos(t)=$\dotfill ou $\cos(t)=$\dotfill  
	\item Si $P\in$ Quadrant \RN{1} et $(\vv{OI}~;~\vv{OP})=t+2k\pi$, alors $\cos t\ldots 0$ et $\sin t \ldots 0$.
\\ Si $P\in$ Quadrant \RN{4} et $(\vv{OI}~;~\vv{OP})=t+2k\pi$, alors $\cos t\ldots 0$ et $\sin t \ldots 0$.
	\item La fonction $\cos$ est de période \dotfill. Donc $\cos(t+\ldots\ldots)=\dotfill$.
	\item La fonction $\sin$ est \dotfill. Pour tout $t\in\mathbb{R}$ on a $\sin(-t)=$\dotfill   
	 \item  $\sin(\pi)=$\dotfill$\sin\tfrac{\pi}{4}= $\dotfill $\sin\tfrac{\pi}{3}=$\dotfill 
	 $\cos(0)=$\dotfill$\cos(-\tfrac{\pi}{4})= $\dotfill $\cos\tfrac{\pi}{3}=$\dotfill  
	  \\ Si   $\cos(t)=-\tfrac{1}{2}$ et $\sin(t)>0$, alors la mesure principale de $t$ est \dotfill  
	 \\ Si $\sin(t)=\tfrac{1}{2}$, alors la mesure principale de $t$ est \dotfill ou \dotfill 
\end{enumerate}
\end{spacing} \vspace{-0.5em}
\end{exercise}  
\begin{exercise} Complétez : 
\vspace{-1em}
\begin{spacing}{1.8}
	\begin{multicols}{2} 
\begin{enumerate}[label={\arabic*)},leftmargin=*]
	\item Si $\sin(t)=0.2$ alors $\sin(-t)=$\dotfill 
	\item Si $\sin(t)=0.35$ alors $\sin(\pi-t)=$\dotfill 
	\item Si $\sin(t)=-0.6$ alors $\sin(\pi+t)=$\dotfill 
	\item Si $\cos(t)=0.8$ alors $\cos(-t)=$\dotfill 
	\item Si $\cos(t)=-0.8$ alors $\cos(\pi-t)=$\dotfill 
	\item Si $\cos(t)=0.1$ alors $\cos(3\pi+t)=$\dotfill
	\item Si  $\cos(t-4\pi)=0.3$ alors $\cos(t)=$\dotfill
	\item Si  $\sin(\pi-t)=0.4$ alors $\sin(t)=$\dotfill  
\end{enumerate} 
	\end{multicols}\vspace{-2.25em}
\end{spacing}
\end{exercise}
\begin{exercise} Complétez afin de déterminer 
	\vspace{-0.5em}
	\begin{spacing}{2}
	\begin{enumerate}[label={\alph*)},leftmargin=*] 
		\item   $\sin\big(\tfrac{8\pi}{3}\big)= \sin\big(\ldots\ldots+2\pi\big) = \sin\big(\ldots\ldots\big)=\sin\big(\pi-\ldots\ldots\big) =\ldots \sin\big(\tfrac{\pi}{3}\big) =\dotfill $ 
		\item $\cos\big(\tfrac{7\pi}{6}\big)= \cos\big(\ldots\ldots+\pi\big) =\ldots\cos\big(\ldots\ldots\big)=$\dotfill
		\item $\sin\big(\tfrac{19\pi}{4}\big)=\ldots\sin(\tfrac{3\pi}{4}+\ldots)=\ldots\sin(\tfrac{3\pi}{4})=\sin\big(\pi-\ldots\ldots\ldots\big)=\ldots\sin\big(\ldots\ldots\big)=$\dotfill
		\item $\cos\big(\tfrac{17\pi}{6}\big)= \cos\big(\tfrac{\pi}{6}+\ldots\ldots \big) = $\dotfill
	\end{enumerate}\vspace{-2em}
\end{spacing}  
\end{exercise}
\begin{exercise}[\cancel{\faCalculator}]  Déterminer les images suivantes en détaillant les étapes.
\vspace{-1em}
\begin{spacing}{1.8}
\begin{multicols}{5} 
	\begin{enumerate}[label={\arabic*)},leftmargin=*]
		\item  $\sin \tfrac{5\pi}{3}$
		\item  $\cos \tfrac{11\pi}{3}$
		\item  $\sin \tfrac{11\pi}{4}$ 
		\item  $\sin (25\pi)$ 
		\item  $\cos (-250\pi)$ 
		\item  $\cos\big(-\tfrac{\pi}{3}\big) $
		\item  $\cos\big(-\tfrac{7\pi}{6}\big) $
		\item  $\sin\big(-\tfrac{2\pi}{3}\big) $
		\item  $\sin\big(-\tfrac{3\pi}{4}\big) $
		\item  $\cos\big(-\tfrac{11\pi}{3}\big) $
	\end{enumerate}
\end{multicols}\vspace{-2.5em}
\end{spacing}
\end{exercise}
\begin{example}[\cancel{\faCalculator}] $t$ est situé dans le Quadrant \RN{4}. Déterminer $\sin(t)$ sachant que $\cos(t)=\tfrac{3}{5}$.	
\end{example}
\begin{proof} $\cos^2(t)+\sin^2(t)=1$ donc $\sin^2(t)=1-\cos^2(t)=1-\left(\tfrac{3}{5}\right)^2=\tfrac{16}{25}$, $\sin(t)=\pm \tfrac{4}{5}$.\\
	$t$ est dans le Quadrant \RN{4}, $\sin(t)<0$ et on a $\sin(t)=-\tfrac{4}{5}$.
\end{proof}
\begin{exercise}[\cancel{\faCalculator}]  Dans chaque cas, déterminer l'image demandée :
\vspace{-0.5em}
\begin{spacing}{1.8} 
\begin{enumerate}[label={\arabic*)},leftmargin=*]
	\item  $\sin(t)=-\tfrac{4}{5}$, déterminer $\cos(t)$ sachant que $t$ est dans le Quadrant \RN{4}.
	\item $\cos(t)=-\tfrac{7}{25}$, déterminer $\sin(t)$ sachant que $t$ est dans le Quadrant \RN{3}.
	\item $\sin(t)=-\tfrac{1}{4}$, déterminer $\cos(t)$ sachant que  $\cos(t)<0$.
	\item  $\cos(t)=-\tfrac{1}{3}$, déterminer $\sin(t)$ sachant que $t$ est dans le Quadrant \RN{4}.
	\item $\sin(t)=0.8$, déterminer $\cos(t)$ sachant que $t\in\interval{\tfrac{\pi}{2}}{\pi}$.
\end{enumerate} \vspace{-1.5em}
\end{spacing}
\end{exercise}
\begin{exercise}[\cancel{\faCalculator}]  Trouver dans chaque cas le réel $t$ demandé.
		\vspace{-1em}
	\begin{multicols}{3}
		\begin{enumerate}[label={$(E_{\arabic*})$}, leftmargin=*]
			\item $\cos(x)=\tfrac{\sqrt{2}}{2}$ et $x\in\interval{-\pi}{0}$
			\item $\cos(x)=-\tfrac{1}{2}$ et $x\in\interval{\pi}{2\pi}$
			\item $\sin(x)=1$ et $x\in\interval{\tfrac{\pi}{2}}{\pi}$
			\item $2\sin(x)=1$ et $x\in\interval{\tfrac{\pi}{2}}{\pi}$
			\item $\sin(x)=\tfrac{1}{2}$ et $x\in\interval{-\pi}{0}$
			\item $2\cos(x)=\sqrt{2}$ et $x\in\interval{-\tfrac{\pi}{2}}{0}$
		\end{enumerate} 
	\end{multicols}   
\end{exercise}
\begin{example} Analyser les résolutions dans $\mathbb{R}$ des équations d'inconnue $x$ suivantes :   \\
		\noindent
		\parbox{0.37\linewidth}{
			$\begin{WithArrows}
				&\sqrt{2}\cos(x)+1 =0  \\%\Arrow{isoler $\cos(x)$} \\ 
				\cos(x)&=-\tfrac{1}{\sqrt{2}}=-\tfrac{\sqrt{2}}{2} \\ %\Arrow{$\cos(\tfrac{3\pi}{4})=-\tfrac{\sqrt{2}}{2}$} \\ 
				\cos(x)&=\cos(\tfrac{3\pi}{4})\\
				x&=\tfrac{3\pi}{4}+2k\pi \quad \text{ou}\quad -\tfrac{3\pi}{4}+2k'\pi
			\end{WithArrows}$ 
		}
	\hfill 
		\parbox{0.6\linewidth}{
			$\begin{WithArrows}
				2\sin(x)& =1  \Arrow{isoler $\cos(x)$ ou $\sin(x)$} \\ 
				\sin(x)&=\tfrac{1}{2} \Arrow{$\cos(\tfrac{3\pi}{4})=-\tfrac{\sqrt{2}}{2}$ et $\sin(\tfrac{\pi}{6})=-\tfrac{1}{2}$} \\ 
				\sin(x)&=\sin(\tfrac{\pi}{6})\\
				x=\tfrac{\pi}{6}+2k\pi \quad &\text{ou}\quad \pi-\tfrac{\pi}{6}+2k'\pi
			\end{WithArrows}$ 
		} 
\end{example}
\begin{exercise}[\cancel{\faCalculator}]  Résoudre dans $\mathbb{R}$ les équations suivantes :\hfill 
	\vspace{-1em}
	\begin{multicols}{4}
		\begin{enumerate}[label={$(E_{\arabic*})$}, leftmargin=*]
			\item $\cos(x)=-\tfrac{1}{2}$
			\item $\cos(x)=-1$
			\item $\sin(x)=\tfrac{\sqrt{3}}{2}$
			\item $2\sin(x)+\sqrt{3}=0$
		\end{enumerate} \vspace{-1em}
	\end{multicols}   
\end{exercise}
\begin{exercise}[parité] Pour chaque fonction définie sur $\mathbb{R}$, comparer les expressions de $f(x)$ et $f(-x)$ puis déterminer si la fonction est paire, impaire ou aucun des deux. Vérifiez graphiquement sur la numworks.
	\vspace{-0.5em}
\begin{multicols}{4}
	\begin{enumerate}[ label={$f_{\arabic*}(x)=$},leftmargin=*]
		\item $x^2\sin(x)$
		\item $\sin(x)\cos(x)$
		\item $x^3+\cos(x)$.
		\item $x^2\cos(2x)$
%		\item $\sin(x)+\cos(x)$
%		\item $x\sin^3(x)$
%		\item $\cos(\sin x)$ 
	\end{enumerate}	\vspace{-1em}
\end{multicols}
\end{exercise}
\begin{eBox}
	Les fonctions du type $t\mapsto A\sin(\omega t+\varphi)$ on une période $T=\tfrac{2\pi}{\omega}$ \href{https://www.desmos.com/calculator/aedicceo7x}{c.f. lien}
\end{eBox}
\begin{exercise} Pour chaque fonction définie sur $\mathbb{R}$, comparer les expressions de $f(x+T)$ avec $f(x)$. En déduire si la fonction est périodique et une période. 	\vspace{-0.5em}
\begin{multicols}{3}
	\begin{enumerate}[ label={$f_{\arabic*}(x)=$},leftmargin=*]
		\item $\sin(x)\cos(x)$ avec $T=\pi$
		\item $\cos(5x+3)$ avec $T=\tfrac{2\pi}{5}$
		\item $2\sin(3x-1)$  avec $T=\tfrac{2\pi}{3}$ 
	\end{enumerate}	\vspace{-0.5em}
\end{multicols}
\end{exercise}
\begin{exercise}[calcul de $\cos\tfrac{\pi}{3}$ et $\sin\tfrac{\pi}{3}$] Soit un repère orthonormé $(O~;~I,J)$ orienté, et le cercle trigonométrique et les point  $P(x~;~y)$ correspondant  à $\tfrac{\pi}{3}$ et $Q$ celui pour $\tfrac{2\pi}{3}$. Par symétrie on peut écrire $Q(-x~;~y)$.
\noindent
\parbox{0.25\linewidth}{\centering 
\begin{tikzpicture}[scale=2]
	\tkzSetUpPoint[size=3, fill=black]%,shape=cross out
	\tkzfctset{tan style/.style={->,>=latex, color=red, line width=1.2pt}}   
	\tkzSetUpLine[line width= 1.2pt, add=.2 and .2]  
	\tkzInit[xmin=-1.0,xmax=1.0,ymin=-0.0,ymax=1.0,xstep=1.0, ystep=1.0];   
	\tkzDefPoint(0,0){O}  
	\tkzDefPoint(1,0){I}   
	\tkzDefPoint(-1,0){I'}  
	\tkzDefPoint(0,1){J}  
	\tkzDrawX[ font=\small, line width=1.2pt, label={}  ,below left, left space=0.1,right space=0.1] 
	\tkzDrawY[ font=\small, line width=1.2pt , label={} ,below left, down space=0.1,up space=0.1]   
	\tkzDrawArc[line width=2pt,color=black](O,I)(I')
	\foreach \x/\xtext in {  
		0/I,
		90/J,
		60/P,
		1200/Q 
	} 
	{
		\tkzDefShiftPoint[O](\x:1){\xtext};
		\tkzDrawPoints(\xtext);    
	} 
	\tkzDefShiftPoint[O](60:1){P};
	\tkzDefShiftPoint[O](120:1){Q};
	\tkzLabelPoint[above ](P){\scriptsize$P(x~;~y)$}
	\tkzLabelPoint[above ](Q){\scriptsize$Q(-x~;~y)$} 
	\tkzDrawSegments[](I,P P,Q Q,I')
 	\tkzMarkSegments[mark=s||](I,P P,Q Q,I')
\end{tikzpicture} 
}
\parbox{0.74\linewidth}{
\begin{enumerate}[leftmargin=*,label=\arabic*)]
	\item Utilisez la formule de la distance, exprimer $PQ^2$ et $IP^2$ en fonction de $x$.
	\item Justifiez que $x$ et $y$ vérifient le système : $\begin{cases}
		4x^2&=(x-1)^2+y^2\\
		x^2+y^2&=1
	\end{cases}$.
	\item Par substitution, montrer que $x$ est solution de $2x^2+x-1=0$
	\item Résoudre et déduire les valeurs admissibles pour $x$ et $y$.
\end{enumerate}	
}
\end{exercise}
%----------------------------------------------------------------------------------------
%	 
%----------------------------------------------------------------------------------------
\newpage 
\loadgeometry{widelayout}  
\pagestyle{scrheadings} % Print headers again  
\KOMAoptions{
	headwidth=textwithmarginpar,
	headsepline=true,
	headsepline= 0.5pt,
	footwidth=textwithmarginpar,
}  
\onehalfspacing 
\subsection{Corrections et éléments de réponses} 

%----------------------------------------------------------------------------------------
%	 
%----------------------------------------------------------------------------------------
\end{document}